% \iffalse meta-comment
% 
% Copyright (C) 2014 by Carsten Burgard
% 
% This file may be distributed and/or modified under the
% conditions of the LaTeX Project Public License, either
% version 1.3 of this license or (at your option) any later
% version. The latest version of this license is in:
% 
%     http://www.latex-project.org/lppl.txt
% 
% and version 1.3 or later is part of all distributions of LaTeX
% version 2005/12/01 or later.
%
% The pictures accompanying this package
%
%       arrow.pdf  arrow.eps
%       box.pdf    box.eps
%       scroll.pdf scroll.eps
%       water.pdf  water.eps
%       fire.pdf   fire.eps
%       earth.pdf  earth.eps
%       air.pdf    air.eps
%       ice.pdf    ice.eps
%       wood.pdf   wood.eps
%       magic.pdf  magic.eps
%       metal.pdf  metal.eps
%
% are property of 
%
%   	Verlag für Fantasy- und Science Fiction-Spiele (VFSF)
%       Elsa Franke
%	Ringstraße 22
%       D-67705 Stelzenberg
%
% and are used and distributed with kind permission thereof.
% 
% \fi
%
% \iffalse
%<package>\NeedsTeXFormat{LaTeX2e}[1999/12/01]
%<package>\ProvidesPackage{midgard}
%<package> [2014/01/26 v1.0 A package for typesetting documents for the german Pen & Paper Roleplay Game 'Midgard']
%<*driver>
\documentclass{ltxdoc}
\usepackage[ngerman,english]{babel}
\usepackage[utf8]{inputenc}
\usepackage[T1]{fontenc}
\usepackage{midgard}
\usepackage{flafter}
\usepackage{float}
\usepackage{placeins}
\usepackage{changepage}
\usepackage[margin=10pt,font=small,labelfont=bf,aboveskip=1em,labelsep=endash]{caption}
\usepackage[a4paper,lmargin=3.5cm, rmargin=3.5cm]{geometry}
\usepackage{array}
\usepackage{boxedminipage}
\usepackage{hyperref}
\newenvironment{widecenter}{%
  \begin{adjustwidth}{0.5\textwidth-0.5\paperwidth+1cm}{0.5\textwidth-0.5\paperwidth+1cm}
  \setlength\LTleft{0pt plus 1fill minus 1fill}%
  \let\LTright\LTleft
  \centering
}{%
  \end{adjustwidth}
}
\newcommand{\footnoteremember}[2]{%
  \footnote{#2}%
  \newcounter{#1}%
  \setcounter{#1}{\value{footnote}}%
}   
\def\generalname{Allgemein}
\GlossaryPrologue{%
  \section*{Änderungen}
}
\IndexPrologue{%
\section*{Index}
In \textit{Kursivschrift} gesetzte Zahlen beziehen sich auf die Seite,
auf welcher der betreffende Eintrag beschrieben
wird. \underline{Unterstrichene} Zahlen beziehen sich auf die
Code-Zeilen der Definition. Zahlen in Normalschrift beziehen sich auf
Code-Zeilen, in denen der betreffende Eintrag verwendet wird.  }
\setlength{\parskip}{3pt}
\setcounter{topnumber}{2}
\setcounter{bottomnumber}{2}
\setcounter{totalnumber}{4}
\setcounter{dbltopnumber}{2}
\renewcommand{\topfraction}{0.9}	
\renewcommand{\bottomfraction}{0.8}
\renewcommand{\textfraction}{0.07}
\newcommand{\footnoterecall}[1]{%
  \footnotemark[\value{#1}]%
}
\EnableCrossrefs
\CodelineIndex
\RecordChanges
\begin{document}
% \OnlyDescription
\DocInput{midgard.dtx}
\end{document}
%</driver> \fi \CheckSum{0} \changes{v1.0}{2012/05/23}{Ursprüngliche
% Version} \GetFileInfo {midgard.sty}
% \makeatletter
% \expandafter\DoNotIndex\expandafter{\string\{}
% \expandafter\DoNotIndex\expandafter{\string\}}
% \DoNotIndex{\#,\$,\%,\&,\@,\\,\^,\_,\~,\,,\.,\=,\>,\@ne,\advance,\begingroup,\catcode,\closein,\closeout,\day,\def,\edef,\else,\empty,\endgroup,\if,\else,\fi,\begin,\end,\leftskip,\let,\detokenize,\expandafter,\linewidth,\newcommand,\newcounter,\newenvironment,\newline,\noindent,\parindent,\parsep,\parskip,\partopsep,\setcounter,\setlength,\textbf,\textit,\textsf,\textsc,\text...,\large,\Large,\small,\footnotesize,\itshape,\bfseries,\par,\textup,\times,\topmargin,\topsep,\topskip,\LARGE,\headsep,\equal,\value,\frac,\selectfont,\setkeys,\savekeys,\ref,\rightskip,\raggedright,\sloppy,\textfloatsep,\texthigh,\vspace,\togglefalse,\toggletrue,\typeout,\usetikzlibrary,\xdef,\newlength,\newmdenv,\nolinebreak,\nopagebreak,\nicefrac,\numdef,\numgdef,\mdfdefinestyle,\makebox,\kill,\includegraphics,\iftoggle,\ifnumgreater,\ifdefstring,\hfill,\fi,\csname,\endcsname,\guillemotleft,\guillemotright,\thinspace,\appto,\addvspace,\AtBeginDocument,\bgroup,\centering,\cleaders,\col@number,\copy,\define@key,\everypar,\ensuremath,\expandnext,\flushleft,\fontsize,\global,\glqq,\grqq,\hfil,\hskip,\Hy@raisedlink,\ifboolexp,\ifcsdef,\ifcsname,\ifdefstrequal,\ifnum,\ifodd,\ifstrempty,\lastbox,\linebreak,\mbox,\newsavebox,\newtoggle,\nobreak,\numexpr,\pagebreak,\parbox,\presetkeys,\protect,\providecommand,\protect,\RequirePackage,\romannumeral,\savebox,\setbox,\the,\unpenalty,\unskip,\unvbox,\vbox,\z@,\zz,\relax,\egroup,\endtabbing,\ifboolexpr,\hspace,\hyphenation,\hypertarget,\hyperlink,\@afterheading,\@afterindentfalse,\@ifpackageloaded,\@ifundefined,\addcontentsline,\hyphenation,\stepcounter,\tabbing}
% \makeatother
% \title{Das \textsf{ midgard } Paket\thanks{Dieses Dokument bezieht
% sich auf \textsf{ midgard }~\fileversion, datiert~\filedate.}}

% \author{ Carsten Burgard \\ \texttt{ carsten.burgard@gmail.com }}

% \maketitle 

% \selectlanguage{ngerman}

% \begin {abstract}

% \textsc{Midgard} ist ein deutsches Pen-\&-Paper Fantasy Rollenspiel,
% das seit 1981 erscheint und damit das erste deutsche
% Fantasy-Rollenspiel überhaupt war.  Dieses \LaTeX{}-Paket versucht,
% das Layout der offiziellen, häufig mit propietären
% Textsatzprogrammen gesetzten \textsc{Midgard}-Publikationen der
% vierten Auflage so gut es geht nachzuahmen. Zu diesem Zweck werden
% einige Kommandos und Umgebungen definiert. Insbesondere wurde
% versucht, ein störungsfreies und vollautomatisches
% Zwei-Spalten-Layout zu erlauben.

% \end{abstract}

% \selectlanguage{english}

% \begin{abstract} \textsc{Midgard} is a
% german pen \& paper roleplay game. This \LaTeX{} package tries to
% resemble the layout of the official publications of this game and
% make it possible to write your own documents in a similar style and
% layout. Different commands and environments are defined to fit this
% purpose. The package and the documentation is -- apart from this
% paragraph -- entirely written in german, as is  the game itself. A
% translation of this package into other languages is not planned,
% since there is little use for this package outside the entirely
% german-speaking fanbase of the game itself. 

% \end {abstract}

% \selectlanguage{ngerman}

% \section{Einleitung} 

% Spielleiter verbringen mitunter viel Zeit mit ihrer
% Kampagnenplanung. Es gibt zwar inzwischen eine große Zahl von
% Abenteuern, sei es im Internet zum freien oder kostenpflichtigen
% Download oder in gedruckter Form, als Hefte oder Teile von
% Quellenbüchern -- doch letzten Endes kommt man nicht umhin, immer
% wieder Zeit aufzuwenden, um eigene Abenteuer oder Regelerweiterungen
% zu entwickeln. Oft jedoch ist die Hemmschwelle groß, die eigenen
% Ideen anderen zugänglich zu machen -- denn das würde bedeuten, dass
% man sie zunächst einigermaßen ordentlich aufschreiben muss. Viel
% Zeit geht dabei verloren, die Werte von neuen Zaubern oder Kreaturen
% abzutippen und in eine einigermaßen ansehnliche, wenn nicht sogar
% mit den offiziellen Schreibweisen konforme Form zu bringen.  Dieses
% Paket versucht, diese Hemmschwelle zu senken, indem es Kommandos und
% Umgebungen zur Verfügung stellt, die es erlauben, auf einfach und --
% hoffentlich -- intuitive Weise ein Layout zu erzielen, das dem der
% offiziellen Veröffentlichungen so nahe wie möglich kommt. Ich möchte
% jedoch nicht verhehlen, dass ich bei einigen Details auch bewusst
% vom offiziellen Layout abweiche, um das Layout unter Verwendung der
% Kapazitäten von \LaTeX{} noch ansehnlicher zu machen. Die folgende
% Liste enthält die von mir bewusst abweichend implementierten
% Details.

% \begin{itemize}

% \item Alle sonst in GROSSBUCHSTABEN geschriebenen Begriffe, wie
% Verweise auf Quellenbücher und Regelwerke, werden stattdessen in
% \textsc{Kapitälchen} gesetzt, um sie weniger blockig wirken zu
% lassen.

% \item Wenn die Felder \glqq Besonderheiten\grqq{} und \glqq
% Angriff\grqq{} von Kreaturendaten so lang sind, dass sie
% Zeilenumbrüche verursachen, erhalten die nachfolgenden Zeilen einen
% leichten hängenden Einzug.

% \item Die
% Werte von Kreaturen enthalten ein zusätzliches Feld \glqq Aura\grqq,
% das im Anschluss an die Besonderheiten in die gleiche Zeile gedruckt
% wird. 

% \end{itemize}

% \section{Anleitung}

% Neben einigen mehr oder minder trivialen Kommandos besteht eine der
% Hauptleistungen dieses Pakets in der Implementierung der Kommandos
% |\fertigkeit|, |\waffenfertigkeit|, |\zauber|, |\zauberwerkstatt|
% und |\bestiarium| (bzw.~|\best|), welche ein den offiziellen
% Publikationen nachempfundenes Layout der entsprechenden Datenblöcke
% erlauben. Intern wird hierzu das Paket |xkeyval| verwendet, welches
% die Übergabe einer durch Kommas getrennten Liste von
% Schlüssel-Wert-Paaren im optionalen Argument eines Kommandos
% erlaubt. Ein Aufruf beispielsweise des |\fertigkeit|-Kommandos
% erfolgt durch die Übergabe der Werte sämtlicher Felder im optionalen
% Argument in der Form |\fertigkeit[|\meta{Schlüssel}|=|\meta{Wert}|,|
% \meta{Schlüssel}|=|\meta{Wert}|,|\dots|]{}|. Dies gilt analog für
% die anderen Kommandos zum Textsatz von Datenblöcken. Das Layout und
% die Spezifikationen der Schlüssel-Wert-Paare werden im folgenden
% jeweils detailliert beschrieben.

% \subsection{Allgemeines}

% \DescribeMacro{\midgardabenteuer} Auf der Titelseite von offiziellen
% \textsc{Midgard}-Abenteuern sind oft die Worte \glqq Midgard
% Abenteuer\grqq{} in einer dekorativen Weise angebracht. Einen
% ähnlichen Schriftzug wie den in offiziellen Abenteuern abgedruckten
% liefert das Kommando |\midgardabenteuer|, dessen Ergebnis in
% Abbildung~\ref{fig:midgardabenteuer} dargestellt wird. Je nach
% Kontext sollte es mit |\centering| auf der Seite zentriert
% werden. \DescribeMacro{\midgardquellenbuch} Ein analoges Kommando
% existiert auch für Quellenbücher.

% \begin{figure}[H] \centering \midgardabenteuer \par\bigskip\par \midgardquellenbuch \caption{Schriftzug
% auf der Titelseite von
% Abenteuern und Quellenbüchern\label{fig:midgardabenteuer}} \end{figure}

% \DescribeEnv{beispiel} Gerade im Grundregelwerk \glqq
% \textsc{Midgard} -- Das Fantasy Rollenspiel\grqq{} finden sich immer
% wieder Beispiel-Passagen, welche die konkrete Verwendung bestimmter
% Regeln demonstrieren sollen. Diese sind typischerweise in leicht
% grau hinterlegte, abgerundete Boxen abgedruckt. Dies versucht die
% Umgebung |beispiel| nachzuempfinden. Das Ergebnis ist in
% Abbildung~\ref{fig:beispiel} dargestellt. \DescribeEnv{textbox} Auch
% eckige Textboxen \DescribeEnv{rundbox} und runde Informationsboxen
% sind üblich, wie in Abbildungen~\ref{fig:textbox} und
% \ref{fig:rundbox} dargestellt. Ein optionales Argument erlaubt die
% Eingabe einer Titelzeile für die Box.

% \begin{figure}[H]
% \centering \begin{beispiel} In einer solchen \textit{Beispiel-Box}
% können erklärende Textpassagen abgedruckt werden. Die Breite
% expandiert dabei automatisch auf den gesamten zugänglichen
% Bereich. Soll sie künstlich eingeschränkt werden, kann hierfür die
% Standard-Umgebung |minipage| verwendet
% werden. 
% \end{beispiel}\vspace{\abovecaptionskip}
% \caption{Die
% Umgebung \textsf{beispiel} \label{fig:beispiel}} \end{figure}

% \begin{figure}[H]
% \centering \begin{textbox}[Textbox] In einer solchen \textit{Textbox}
% können zum direkten Vorlesen gedachte Passagen gesetzt werden.
% \end{textbox}\vspace{\abovecaptionskip}
% \caption{Die
% Umgebung \textsf{textbox} \label{fig:textbox}} \end{figure}

% \begin{figure}[H]
% \centering \begin{rundbox}[Rundbox] In einer solchen \textit{Rundbox}
% werden üblicherweise nicht direkt mit dem Haupttext in Verindung stehende
% Hintergrundinformationen abgedruckt.
% \end{rundbox}\vspace{\abovecaptionskip}
% \caption{Die
% Umgebung \textsf{rundbox} \label{fig:rundbox}} \end{figure}

% \DescribeMacro{\geld} Mit |\geld| kann eine beliebige Goldmenge
% korrekt gestückelt ausgegeben werden. Beispielsweise produziert
% |\geld{23456}| die Ausgabe \textit{\geld{23456}}.  Da es eine beliebte
% Hausregel ist, die Umrechnungsfaktoren von Gold- in Silberstücke
% oder von Silber- in Goldstücke anzupassen (beispielsweise auf 1:100
% anstelle des Normalwertes von 1:10) existieren entsprechende
% Funktionalitäten, namentlich |\setzeKupferProSilber| und
% |\setzeSilberProGold|. Außerdem können die Anzeigewerte \textit{GS},
% \textit{SS} und \textit{KS} durch Umdefinition von |\goldAlias|,
% |\silberAlias| und |\kupferAlias| editiert werden.

% \DescribeMacro{\setzeGliederungsEbene} Alle für Datenblöcke
% definiertern Umgebungen fügen dem Inhaltsverzeichnis automatisch
% Einträge hinzu. Die Gliederungsebene, auf welcher diese erscheinen,
% können mit dem Makro
% |\setzeGliederungsEbene{|\meta{Typ}|}{|\meta{Ebene}|}| kontrolliert
% werden.

% \subsection{Das Fantasy Rollenspiel} \begin{figure} \centering

% \begin{minipage}{1.0\textwidth}

% \fertigkeit[name=Schleucheln, 
%             kategorie=Kampf, 
%             ungelernt=0,
%             leitAtt={Gs\,91}, 
%             EW=4,
%             anforderungen={Gw\,81, \itshape Schleichen+12, Meucheln+12},
%             FP=200,
%             G=As,
%             S={\textsc{Käm} a.~As},
%             A={alle anderen}
% ]{}

% Der Kämpfer kann gleichzeitig \textit{Schleichen} und
% \textit{Meucheln}.

% \bigskip

% \fertigkeit[name=Atmen, EW=20, anforderungen={Wk 01}]{}

%  Jeder Mensch kann atmen. Ein kritischer Fehler führt zum
% Verschlucken.

% \end{minipage}\par 

%\bigskip\hrule\bigskip\par 

%\begin{minipage}{0.9\textwidth} \begin{verbatim}

% \fertigkeit[name=Schleucheln, 
%             kategorie=Kampf, 
%             ungelernt=0,
%             leitAtt={Gs\,91}, 
%             anforderungen={Gw\,81, \itshape Schleichen+12,
%             Meucheln+12}, 
%             FP=200,
%             EW=4, 
%             G=As, 
%             S={\textsc{Käm} a.~As}, 
%             A={alle anderen}
% ]{} 

% Der Kämpfer kann gleichzeitig \textit{Schleichen} und
% \textit{Meucheln}.  

% \bigskip 

% \fertigkeit[name=Atmen, 
%             EW=20,
%             anforderungen={Wk 01}
% ]{}
 
% Jeder Mensch kann atmen. Ein kritischer Fehler führt zum
% Verschlucken.

% \end{verbatim} 
% \end{minipage} 
 
% \caption{Ein dem \textsc{Dfr4} nachempfundenes Layout für Allgemeine Fertigkeiten\label{fig:fertigkeit}} 
% \end{figure}

% \DescribeMacro{\fertigkeit} Mit dem Kommando
% |\fertigkeit[|\meta{Schlüssel}|=|\meta{Wert}|]{}| kann ein dem
% \textsc{Dfr4} nachempfundenes Layout von allgemeinen Fähigkeiten
% erzeugt werden. Die Benutzung und das Ergebnis wird durch
% Abbildung~\ref{fig:fertigkeit} an zwei (nicht ganz ernstzunehmenden)
% Beispielen illustriert. Das erste stellt eine reguläre allgemeine,
% das zweite eine angeborene Fertigkeit ohne Lernmechanismus
% dar. Letzteres wurde erzielt, indem der Wert |FP| auf \textit{hide}
% gesetzt bzw.~auf seinem Standard-Wert belassen wurde. Die
% Tabelle~\ref{tbl:fertigkeit} enthält eine vollständige Liste aller
% implementierten Schlüssel-Wert-Paare. 

% \begin{widecenter} \setkeys{fertigkeit}{}{ \begin{longtable}{>{\sf}l
% | >{\itshape}c | p{12cm}} \caption{Vollständige Liste der Schlüssel
% und Felder des
% \textsf{fertigkeit}-Kommandos\label{tbl:fertigkeit}}\\ \rm\bfseries
% Schlüssel & default & \bfseries Wirkung \tabularnewline\hline name &
% \fertigkeitName & Der Name der Fertigkeit wird durch
% \textbf{Fettdruck} hervorgehoben in der Titelzeile dargestellt. Eine
% Besonderheit besteht, wenn |FP| nicht gesetzt ist -- dann ist die
% Fähigkeit nicht lernbar und hat daher einen fixen Erfolgswert, der
% dann unmittelbar auf den Namen folgend dargestellt
% wird.\tabularnewline entry & \fertigkeitEntry & Eintrag dieser
% Fähigkeit im Inhaltsverzeichnis.  Bei \textit{\empty} wird der mit
% \meta{name} gesetzte Name der Fertigkeit verwendet. \tabularnewline
% kategorie & \fertigkeitKategorie & Die Kategorie der Fähigkeit, etwa
% \glqq Sozial\grqq{} oder \glqq Bewegung\grqq{}. Jeder Wert außer
% \textit{\empty} wird in der Titelzeile hinter dem Namen der Fähigkeit
% in Klammern abgedruckt. \tabularnewline ungelernt &
% \fertigkeitUngelernt & Der ungelernte Erfolgswert für diese
% Fertigkeit. Jeder Wert außer \textit{\empty} wird als Erfolgswert
% interpretiert und in der Form \glqq
% ungelernt+(\meta{ungelernt})\grqq{} am Ende der Titelzeile
% abgedruckt. Der Löwenanteil des folgenden Datenblocks wird dann
% ignoriert und nur das Leitattribut und die anderen Anforderungen
% werden noch in der Titelzeile dargestellt, falls diese nicht
% \textit{\empty} sind. \tabularnewline leitatt & \fertigkeitLeitAtt &
% Das Leitattribut dieser Fertigkeit. Jeder Wert außer \textit{\empty}
% wird in der unmittelbar auf den Titel folgenden Zeile in
% \textbf{Fettdruck} an den Anfang gestellt.\tabularnewline
% anforderungen & \fertigkeitAnforderungen & Die zum Lernen dieser
% Fertigkeit benötigten Voraussetzungen, typischerweise
% Mindestattribute. Jeder Wert außer \textit{\empty} wird in der
% unmittelbar auf den Titel folgenden Zeile an den Anfang gestellt,
% hinter das Leitattribut. \tabularnewline EW & \fertigkeitEW & Der
% Wert |EW| wird als Erfolgswert interpretiert und wörtlich in der
% Form \glqq Erfolgswert+\meta{EW}\grqq{} in der Mitte der zweiten
% Zeile des Datenblocks, hinter den Voraussetzungen,
% abgedruckt.\tabularnewline spielbeginn & \fertigkeitEW & Der Wert
% |spielbeginn| wird als Erfolgswert nach Erlernen zu Spielbeginn
% interpretiert und als erster von zwei Werten wörtlich innerhalb der
% Klammer (+\meta{spielbeginn}/\dots) am Ende der zweiten Zeile des
% Datenblocks abgedruckt. \tabularnewline maximum & \fertigkeitEW &
% Der Wert |maximum| wird als maximal erreichbarer Erfolgswert
% interpretiert und als zweiter der beiden Werte wörtlich innerhalb
% der Klammer (\dots/+\meta{maximum}) am Ende der zweiten Zeile des
% Datenblocks abgedruckt. \tabularnewline FP & \fertigkeitFP & Die zum
% Erlernen der Grundkenntnisse in dieser Fertigkeit benötigten
% Erfahrungspunkte. Jede ganze Zahl wird vollständig als
% Standardkosten, halbiert und abgerundet und als Grundkosten sowie
% verdoppelt als Ausnahmekosten abgedruckt. Angabe von \textit{hide}
% führt dazu, dass die Zeile mit den Lernkosten wie auch der Großteil
% des gesamten Datenblocks nicht abgedruckt wird. Dies ist in der
% unteren Hälfte des Beispiels in Abbildung~\ref{fig:fertigkeit} zu
% sehen. Das Verhalten für Angaben außer ganzen Zahlen und
% \textit{hide} ist nicht definiert. \tabularnewline G & \fertigkeitG
% & Diejenigen Abenteurertypen, die diese Fertigkeit zu Grundkosten
% lernen können.\tabularnewline S & \fertigkeitS & Diejenigen
% Abenteurertypen, die diese Fertigkeit zu Standardkosten lernen
% können.\tabularnewline A & \fertigkeitA & Diejenigen
% Abenteurertypen, die diese Fertigkeit zu Ausnahmekosten lernen
% können.
% \end{longtable}}
% \end{widecenter} 
% \begin{figure}
% \centering
% \begin{minipage}{0.6\textwidth} 
% \waffenfertigkeit[
% name=Spazierstockwerfen, schaden=1W6,
% reichweite={5-10/11-20/21-30\,m}, 
% voraussetzungen={St\,31, Gw\,11},
% grundkenntnisse=Alltagsgegenstände, 
% beispiele={Handtasche, Stuhl}, FP=500, 
% G=ältere Herren, S=Wanderer, A=niemand
% ]{} \par
% Beschreibungstext der Waffenfähigkeit 
% \end{minipage}
% \par
% \bigskip
% \hrule
% \bigskip
% \par 
% \begin{minipage}{0.6\textwidth} 
% \begin{verbatim}
% \waffenfertigkeit[
%   name=Spazierstockwerfen, schaden=1W6,
%   reichweite={5-10/11-20/21-30\,m}, 
%   voraussetzungen={St\,31, Gw\,11},
%   grundkenntnisse=Alltagsgegenstände, 
%   beispiele={Handtasche, Stuhl}, FP=500, 
%   G=ältere Herren, S=Wanderer, A=niemand
% ]{} 
% Beschreibungstext der Waffenfähigkeit
% \end{verbatim}
% \end{minipage}
% \caption{Ein dem \textsc{Dfr4} nachempfundenes Layout für
% Waffenfertigkeiten\label{fig:waffe}} 
% \end{figure}

% \DescribeMacro{\waffenfertigkeit} Mit dem Kommando
% |\waffenfertigkeit[|\meta{Schlüssel}|=|\meta{Wert}|,| \dots|]{}|
% kann ein dem \textsc{Dfr4} nachempfundener Textsatz für
% Waffenfertigkeiten erzeugt werden. Das Ergebnis ist im
% Abbildung~\ref{fig:waffe} dargestellt. Die vollständige Liste mit
% Schlüsseln ist in Tabelle \ref{tbl:waffe} abgedruckt. Die Spalte
% \textit{default} beinhaltet dabei denjenigen Wert, den das
% entsprechende Feld annimmt, wenn der Wert nicht aktiv gesetzt wird.
% \DescribeMacro{\weapon} Aus Bequemlichkeits- und
% Kompatibilitätsgründen wird auch das Makro |\weapon| angeboten, das
% sich in jeder Hinsicht aber genauso wie |\waffenfertigkeit|
% verhält.  \begin{widecenter}
% \setkeys{weapon}{}{ \begin{longtable}{>{\sf}l | >{\itshape}c |
% p{12cm}} \caption{Vollständige Liste der Schlüssel und Felder des
% \textsf{waffenfertigkeit}-Kommandos\label{tbl:waffe}}\\ \rm\bfseries
% Schlüssel & default & \bfseries Wirkung \tabularnewline\hline name &
% \weaponName & Der Name der Waffenfertigkeit wird durch
% \textbf{Fettdruck} hervorgehoben in der Titelzeile
% dargestellt.\tabularnewline entry & \weaponEntry & Eintrag dieser
% Waffenfähigkeit im Inhaltsverzeichnis.  Bei \textit{\empty} wird der
% mit \meta{name} gesetzte Name der Waffenfertigkeit verwendet. So ist
% es möglich, z.\,B.~die \glqq\textbf{Schwere Armbrust}\grqq{} als
% \glqq\textbf{Armbrust, schwere}\grqq{} im Inhaltsverzeichnis zu
% führen.  \tabularnewline schaden & \weaponDamage & Der Schaden, den
% die Waffe angerichtet. Er wird in Klammern hinter dem Namen der
% Waffe abgedruckt. \tabularnewline reichweite & \weaponRange & Die
% Reichweite wird an den rechten Rand der Titelteile abgedruckt. Der
% \textit{default}-Eintrag ist die leere Zeichenkette
% \glqq\grqq{}. \tabularnewline voraussetzungen & \weaponPremise & Die
% zum Lernen dieser Waffenfähigkeit benötigten Voraussetzungen,
% typischerweise Mindestattribute. Sie werden in der unmittelbar auf
% den Titel folgenden Zeile an den Anfang gestellt. \tabularnewline
% grundkenntnisse & \weaponBasics & Zum Verwenden der Waffe benötigte
% Grundkenntnisse. Sie werden \textbf{fettgedruckt} in der zweiten
% Zeile nach dem Titel dargestellt. \tabularnewline FP & \weaponFP &
% Die zum Erlernen der Grundkenntnisse in dieser Waffenfertigkeit
% benötigten Erfahrungspunkte. Jede ganze Zahl wird vollständig als
% Standardkosten, halbiert und abgerundet und als Grundkosten sowie
% verdoppelt als Ausnahmekosten abgedruckt. Angabe von \textit{hide}
% führt dazu, dass die Zeile mit den Lernkosten wie auch die
% Schwierigkeit der Waffe nicht abgedruckt wird. Das Verhalten für
% Angaben außer ganzen Zahlen und \textit{hide} ist nicht
% definiert. \tabularnewline schwierigkeit & \weaponDifficulty & Die
% Schwierigkeit dieser Waffenfähigkeit, die sich in den Kosten zum
% Steigern des Erfolgswertes niederschlägt. Sie wird in Normaldruck,
% mit dem vorangestellten Wort \glqq Schwierigkeit\grqq{} in
% \textit{kursiv}, in der Zeile vor den Lernkosten abgedruckt -- aber
% nur dann, wenn |FP| nicht auf \textit{hide} gesetzt
% ist. \tabularnewline G & \weaponG & Diejenigen Abenteurertypen, die
% diese Waffenfertigkeit zu Grundkosten lernen können.\tabularnewline
% S & \weaponS & Diejenigen Abenteurertypen, die diese
% Waffenfertigkeit zu Standardkosten lernen können.\tabularnewline A &
% \weaponA & Diejenigen Abenteurertypen, die diese Waffenfertigkeit zu
% Ausnahmekosten lernen
% können.  \end{longtable}} \end{widecenter} 


% \subsection{Arkanum}

% \begin{table}
% \centering
% \begin{tabular}{c l}
% \imgMagisch &|\imgMagisch |\tabularnewline
% \imgModus   &|\imgModus   |\tabularnewline
% \imgObjectum&|\imgObjectum|\tabularnewline
% \imgFeuer   &|\imgFeuer   |\tabularnewline
% \imgWasser  &|\imgWasser  |\tabularnewline
% \imgErde    &|\imgErde    |\tabularnewline
% \imgLuft    &|\imgLuft    |\tabularnewline
% \imgEis     &|\imgEis     |\tabularnewline
% \imgHolz    &|\imgHolz    |\tabularnewline
% \imgMetall  &|\imgMetall  |\tabularnewline
% \imgMagan   &|\imgMagan   |
% \end{tabular}
% \caption{Besondere Symbole des Arkanums und ihre zugehörigen \LaTeX-Kommandos.\label{tab:symbols}}
% \end{table}

% \begin{figure}
% \begin{widecenter}
% \begin{minipage}{8cm}
% \zauber[name=Neuer Zauber, 
%   scroll=false,th=S, 
%   anmerkung={nur zu Demonstrationszwecken}, 
%   stufe=1,art=Gestenzauber,
%   prozess=Veränden,agens=Holz,reagens=Erde,
%   AP=1, Zd=1\,sec, Rw=30\,m,
%   Wz=Geist, Wb=Zauberer,
%   Wd=\textinfty, Ur=druidisch, FP=100,
%   S={Hl, Hx, \textsc{Pri} a.~T},
%   G={Dr, Th},
%   A=niemand,
%   material={Haare eines Hundes},
%   matpreis={1~KS},
%   matannotation={wird nur bei
%   krit.~Fehler verbraucht}
% ]{}
% Beschreibungstext des Zaubers 
% \end{minipage}
% \hspace{0.5cm}\vrule\hspace{0.5cm} 
% \begin{minipage}{6cm} 
% \begin{verbatim}
% \zauber[name=Neuer Zauber, 
%   scroll=false,th=S, 
%   anmerkung={nur zu Demonstrationszwecken}, 
%   stufe=1,art=Gestenzauber,
%   prozess=Veränden,agens=Holz,reagens=Erde,
%   AP=1, Zd=1\,sec, Rw=30\,m,
%   Wz=Geist, Wb=Zauberer,
%   Wd=\textinfty, Ur=druidisch, FP=100,
%   S={Hl, Hx, \textsc{Pri} a.~T},
%   G={Dr, Th},
%   A=niemand,
%   material={Haare eines Hundes},
%   matpreis={1~KS},
%   matannotation={wird nur bei
%   krit.~Fehler verbraucht}
% ]{}
% Beschreibungstext des Zaubers 
% \end{verbatim}
%  \end{minipage}
%  \end{widecenter}
%  \caption{Ein dem Arkanum nachempfundenes Layout für Zauber\label{fig:zauber}} 
% \end{figure} 

% \DescribeMacro{\img...} Im Arkanum wird eine Reihe von Sonderzeichen
% verwendet, welche im Rahmen dieses Paketes zur Verfügung gestellt
% werden sollen. Die Tabelle \ref{tab:symbols} beinhaltet die Kommandos und
% zugehörigen Symbole.

% \DescribeMacro{\zauber} Mit dem Kommando
% |\zauber[|\meta{Schlüssel}|=|\meta{Wert}|,| \dots|]{}| kann ein dem
% Layout des Arkanums nachempfundener Textsatz der Werte eines Zaubers
% erzeugt werden. Das Ergebnis ist im Abbildung~\ref{fig:zauber}
% dargestellt. Die vollständige Liste mit Schlüsseln ist in Tabelle
% \ref{tbl:zauber} abgedruckt. Die Spalte \textit{default} beinhaltet
% dabei denjenigen Wert, den das entsprechende Feld annimmt, wenn der
% Wert nicht aktiv gesetzt wird.

% \DescribeMacro{\zaubervariante} Das Kommando
% |\zaubervariante{|\meta{Variante}|}| soll benutzt werden um am Ende
% einer Zauberbeschreibung die Variationen des Zaubers für
% verschiedene Zaubererklassen genauer zu beschreiben, etwa
% |\zaubervariante{Thaumaturgie}| für einen Abschnitt, der mit
% \glqq\textbf{\textsc{Thaumaturgie}:} \dots\grqq{} beginnt.

% \DescribeMacro{\zaubermittel} Mit dem Kommando
% |\zaubermittel[|\meta{Schlüssel}|=|\meta{Wert}|,| \dots|]{}| kann
% analog ein dem Layout des Arkanums nachempfundener Textsatz der
% Werte eines Zaubermittels erzeugt werden.

% \DescribeMacro{\zauberwerkstatt} Aus Gründen der Bequemlichkeit und
% Kompatibilität definieren wir zusätzlich das Kommando
% |\zauberwerkstatt[|\meta{Schlüssel}|=|\meta{Wert}|,| \dots|]{}|,
% welches sich jedoch genauso verhält wie
% |\zaubermittel[|\meta{Schlüssel}|=|\meta{Wert}|,| \dots|]{}|.

% \FloatBarrier \begin{widecenter}
% \setkeys{zauber}{}{ \begin{longtable}{>{\sf}l | >{\itshape}c |
% p{12cm}} \caption{Vollständige Liste der Schlüssel und Felder des
% \textsf{zauber}-Kommandos\label{tbl:zauber}}\\ \rm\bfseries
% Schlüssel & default & \bfseries Wirkung \tabularnewline\hline name &
% \zaubername & Der Name des Zaubers wird mit \textbf{Fettdruck} und
% leicht größerer Schriftart hervorgehoben in der Titelzeile
% dargestellt.  \tabularnewline entry & \zauberentry & Eintrag dieses
% Zaubers im Inhaltsverzeichnis. Bei \textit{\empty} wird der mit
% \meta{name} gesetzte Name des Zaubers verwendet. So ist es möglich,
% z.\,B.~die \glqq\textbf{Blaue Bannsphäre}\grqq{} als
% \glqq\textbf{Bannsphäre, Blaue}\grqq{} im Inhaltsverzeichnis zu
% führen.  \tabularnewline scroll & \zauberscroll & Bool'sche Angabe,
% ob dieser Zauber von Schriftrolle gelernt werden kann. Jeder andere
% Wert außer \textit{true} führt dazu, dass das Schriftrollensymbol
% \imgMagisch{} in der Kopfzeile
% abgedruckt wird.\tabularnewline th & \zauberth & Information über
% die thaumaturgische Variante des Zaubers, wird typischerweise auf
% \textit{R} oder \textit{S} für Runenstäbe bzw.~Siegelzauber
% gesetzt. Tatsächlich wird aber jede Angabe außer \textit{\empty} in
% der Kopfzeile abgedruckt. Länge Angaben sollten vermieden werden, da
% diese zu unerwünschten Zeilenüberläufen führen können!
% \tabularnewline anmerkung & \zauberAnmerkung & Anmerkungen zum
% Zauber, etwa zur Quelle, aus welcher der Zauber stammt. Jeder Wert
% außer \textit{\empty} wird unterhalb der Titelzeile klein,
% linksbündig, fett und kursiv agbedruckt. \tabularnewline art &
% \zauberart & Art des Zaubers, z.\,B.~\textbf{Gedankenzauber},
% \textbf{Wortzauber} oder \textbf{Gestenzauber}. Jede Angabe wird vor
% dem Wort \glqq der\grqq{} in der Zeile \glqq Gestenzauber der Stufe
% 1\grqq{} im Beispiel (Abbildung~\ref{fig:zauber})
% abgedruckt. \tabularnewline stufe & \zauberstufe & Stufe des
% Zaubers. Jede Angabe außer \textit{GM} wird mit dem vorangestellten
% Wort \glqq Stufe\grqq{} wörtlich abgedruckt, während \textit{GM} zu
% \glqq Großen Magie\grqq{} (z.\,B.~in \glqq Gestenzauber der
% \textbf{Großen Magie}\grqq{}) expandiert wird \tabularnewline
% material & \zaubermaterial & Jedes Zaubermaterial außer
% \textit{\empty} wird in kursiv unterhalb der Zeile mit Art und Stufe
% des Zaubers abgedruckt. Ist das \meta{material} \textit{\empty},
% werden die Angaben \meta{matpreis} und \meta{matannotation}
% ignoriert. \tabularnewline matpreis & \zaubermatpreis & Der
% Materialpreis. Jede Angabe außer \textit{\empty} wird hinter dem
% Zaubermaterial in Normalschrift und geklammert
% abgedruckt. \tabularnewline matannotation & \zaubermatannotation &
% Weitere Anmerkungen zum Zaubermaterial, beispielsweise
% \glqq\textbf{wird nicht verbraucht}\grqq{} oder \glqq\textbf{wird
% nur bei kritischem Fehler verbraucht}\grqq{}. Jede Angabe außer
% \textit{\empty} wird am Ende der \glqq Zaubermaterial-Zeile\grqq{},
% mit einem Semikolon abgetrennt, abgedruckt. Bei Zaubern, die mehrere
% Zaubermaterialien mit separaten Preisangaben haben, empfiehlt es
% sich, den entsprechenden Passus unmittelbar in \meta{material} zu
% schreiben und die Felder \meta{matkosten} und \meta{matannotation}
% zu ignorieren. \tabularnewline feld & \beschwoerungsfeld & Bei
% Beschwörungen verwendetes Feld. Jede Angabe außer \textit{\empty} wird
% unterhalb der Material-Zeile in Kursiv abgedruckt abgedruckt, wobei
% die Bezeichnung \textit{Köder} in Fettdruck vornangestellt wird. Ist
% er Köder \textit{\empty}, wird die Angabe zum \meta{feldpreis}
% ignoriert. \tabularnewline feldpreis & \beschwoerungsfeldpreis & Der
% Preis des Beschwörungsfelds.  Jede Angabe außer \textit{\empty} wird
% hinter dem Beschwörungsfeld in Normalschrift und geklammert
% abgedruckt. \tabularnewline koeder & \beschwoerungskoeder & Bei
% Beschwörungen verwendete Köder. Jede Angabe außer \textit{\empty} wird
% unterhalb der Material-Zeile in Normalschrift abgedruckt. Ist er
% Köder \textit{\empty}, wird die Angabe zum \meta{koederpreis}
% ignoriert. \tabularnewline koederpreis & \beschwoerungskoederpreis &
% Der Preis des Köders.  Jede Angabe außer \textit{\empty} wird hinter
% dem Köder in Normalschrift und geklammert
% abgedruckt. \tabularnewline \tabularnewline prozess & \zauberprozess
% & Der Prozess-Anteil der Zauberformel. Jede Angabe wird wörtlich vor
% dem Symbol \imgModus{} in der
% Zauberformel abgedruckt.\tabularnewline agens & \zauberagens & Der
% Agens-Anteil der Zauberformel. Jede Angabe wird wörtlich zwischen
% den Symbol \imgModus{} und
% \imgObjectum{} in der Zauberformel
% abgedruckt.\tabularnewline reagens & \zauberreagens & Der
% Regens-Anteil der Zauberformel. Jede Angabe wird wörtlich zwischen
% hinter dem Symbol \imgObjectum{} in der
% Zauberformel abgedruckt.\tabularnewline AP & \zauberAP & Die
% AP-Kosten des Zaubers. Jede Angabe außer \textit{hide} wird wörtlich
% abgedruckt, während \textit{hide} dazu führt, dass die Zeile nicht
% gedruckt wird.\tabularnewline GZ & \zauberGZ & Die Gefährdungszahl
% des Zaubers (nur bei Zaubern der finsteren Magie, siehe
% \textit{Hexenzauber \& Druidenkraft}). Jede Angabe außer
% \textit{hide} wird wörtlich abgedruckt, während \textit{hide} dazu
% führt, dass die Zeile nicht gedruckt wird.\tabularnewline Zd &
% \zauberZd & Die Zauberdauer. Jede Angabe außer \textit{hide} wird
% wörtlich abgedruckt, während \textit{hide} dazu führt, dass die
% Zeile nicht gedruckt wird.\tabularnewline Rw & \zauberRw & Die
% Reichweite des Zaubers. Jede Angabe außer \textit{hide} wird
% wörtlich abgedruckt, während \textit{hide} dazu führt, dass die
% Zeile nicht gedruckt wird.\tabularnewline Wz & \zauberWz & Das
% Wirkungsziel des Zaubers. Jede Angabe außer \textit{hide} wird
% wörtlich abgedruckt, während \textit{hide} dazu führt, dass die
% Zeile nicht gedruckt wird.\tabularnewline Wb & \zauberWb & Der
% Wirkungsbereich des Zaubers. Jede Angabe außer \textit{hide} wird
% wörtlich abgedruckt, während \textit{hide} dazu führt, dass die
% Zeile nicht gedruckt wird.\tabularnewline Wd & \zauberWd & Die
% Wirkungsdauer des Zaubers. Jede Angabe außer \textit{hide} wird
% wörtlich abgedruckt, während \textit{hide} dazu führt, dass die
% Zeile nicht gedruckt wird.\tabularnewline Ur & \zauberUr & Der
% Ursprung des Zaubers, etwa \textbf{dämonisch} oder
% \textbf{elementar}. Jede Angabe außer \textit{hide} wird wörtlich
% abgedruckt, während \textit{hide} dazu führt, dass die Zeile nicht
% gedruckt wird.\tabularnewline FP & \zauberFP & Die FP-Kosten des
% Zaubers. Jede ganze Zahl wird vollständig als Standardkosten,
% halbiert und abgerundet und als Grundkosten sowie verfünffacht als
% Ausnahmekosten abgedruckt. Angabe von \textit{hide} führt dazu, dass
% die Zeile mit den Lernkosten nicht abgedruckt wird. Das Verhalten
% für Angaben außer ganzen Zahlen und \textit{hide} ist nicht
% definiert. \tabularnewline G & \zauberG & Diejenigen
% Zaubererklassen, die den Zauber zu Grundkosten lernen
% können.\tabularnewline S & \zauberS & Diejenigen Zaubererklassen,
% die den Zauber zu Standardkosten lernen können.\tabularnewline A &
% \zauberA & Diejenigen Zaubererklassen, die den Zauber zu
% Ausnahmekosten lernen
% können.  \end{longtable}} \end{widecenter}

% \subsection{Bestiarium}
% \begin{figure}
% \begin{verbatim}
% \bestiarium[ name=Kreatur, mag=true, typ=Humanoid, grad=3, In=m30,
% LP=3W6, AP=3W6+2, MW=18, EP=5, Gw=80, St=70, B=24, RK=KR, 
% abwehr=12, resistenz=13/15/11,
% angriff={Keule+8 (1W6+1) -- Raufen+8 (1W6-2)},
% zauberEW=12, zauber={Bannen von Dunkelheit, Schlaf -- Niessalz},
% zauberArtigEW=18, zauberArtig={Heilen von Wunden \textup{(jede Runde auf sich selbst)}}, 
% bes=kann nur mit magischen Waffen verletzt werden, aura=dämonisch, vorkommen=überall]{}
% \end{verbatim}
% \hrule
% \smallskip
% \hfill
% \begin{minipage}{0.45\textwidth}
% \bestiarium[ name=Kreatur, mag=true, typ=Humanoid, grad=3, In=m30,
% LP=3W6, AP=3W6+2, MW=18, EP=5, Gw=80, St=70, B=24, RK=KR, 
% abwehr=12, resistenz=13/15/11,
% angriff={Keule+8 (1W6+1) -- Raufen+8 (1W6-2)},
% zauberEW=12, zauber={Bannen von Dunkelheit, Schlaf -- Niessalz},
% zauberArtigEW=18, zauberArtig={Heilen von Wunden \textup{(jede Runde auf sich selbst)}}, 
% bes=kann nur mit magischen Waffen verletzt werden, aura=dämonisch, vorkommen=überall]{}
% \end{minipage}\hfill 
% \begin{minipage}{0.45\textwidth}
% \best[ name=Kreatur, mag=true, typ=Humanoid, grad=3, In=m30,
% LP=3W6, AP=3W6+2, MW=18, EP=5, Gw=80, St=70, B=24, RK=KR, 
% abwehr=12, resistenz=13/15/11,
% angriff={Keule+8 (1W6+1) -- Raufen+8 (1W6-2)},
% zauberEW=12, zauber={Bannen von Dunkelheit, Schlaf -- Niessalz},
% zauberArtigEW=18, zauberArtig={Heilen von Wunden \textup{(jede Runde auf sich selbst)}}, 
% bes=kann nur mit magischen Waffen verletzt werden, aura=dämonisch, vorkommen=überall]{}
%\end{minipage}

% \hfill~ \caption{Ein dem Bestiarium nachempfundenes Layout für
% Kreaturendaten\label{fig:bestiarium}} \end{figure}

% \DescribeMacro{\bestiarium} Mit dem Kommando |\bestiarium| kann ein
% an das Bestiarium angelehntes Layout für die Spieldaten von
% Kreaturen -- komplett mit Box -- erzeugt werden. Ähnlich wie schon
% beim Kommando |zauber| werden die Werte der einzelnen Felder dabei
% als Schlüssel-Wert-Paare übergeben. Die Verwendung und das Ergebnis
% sind in Abbildung~\ref{fig:bestiarium} (links) dargestellt. Die
% vollständige Liste aller Schlüssel und Felder findet sich in
% Tabelle~\ref{tbl:bestiarium}. Da die Box unter Verwendung des
% |mdframed|-Pakets erzeugt wird, sind Seiten- und Spaltenumbrüche
% innerhalb der Box grundsätzlich erlaubt. Falls dies nicht erwünscht
% ist, kann die folgende Zeile in der Präambel verwendet werden, um
% Seitenumbrüche innerhalb der |\bestiarum|-Boxen global zu
% unterbinden: |\mdfapptodefinestyle{bestbox}{nobreak}|.

% \DescribeMacro{\best} Mit |\best| kann die aus einigen Abenteuern
% und Zauberbeschreibungen bekannte Kurzform von Kreaturendaten
% erzeugt werden, die ohne Box auskommt. Die Verwendung ist identisch
% mit der des Kommandos |\bestiarium|, das Ergebnis ist ebenfalls sind
% in Abbildung~\ref{fig:bestiarium} (rechts) dargestellt. Die
% vollständige Liste aller Schlüssel und Felder findet sich in
% Tabelle~\ref{tbl:bestiarium}.  \begin{widecenter}
% \setkeys{best}{}{% \begin{longtable}{>{\sf}l | >{\itshape}c |
% p{12cm}} \caption{Vollständige Liste der Schlüssel und Felder des
% \textsf{zauber}-Kommandos\label{tbl:bestiarium}}\\ \rm\bfseries
% Schlüssel & default & \bfseries Wirkung\tabularnewline\hline name &
% \bestName & Der Name des Wesens \tabularnewline mag & \bestMag &
% Falls |mag| \textit{true} ist, wird das Schriftrollensymbol
% \imgMagisch{} hinter nem Namen
% angezeigt.  \tabularnewline addendum & \bestAdd & Bei einigen Wesen
% gibt es einen Namenszusatz, welcher nötigenfalls \textbf{nach} dem
% gegebenenfalls vorhandenen Schriftrollensymbol dargestellt werden
% soll, bei einem Drachen etwa der Zusatz \glqq uralter\grqq{} (für
% die Ausgabe \textbf{Drache
% }\imgMagisch\textbf{, uralter}). Jede Angabe außer \textit{\empty} wird deshalb in Fettdruck und mit
% Komma abgetrennt hinter dem Namen und dem eventuell vorhandenen
% Schriftrollensymbol ausgegeben.  \tabularnewline grad & \bestGrad &
% Der Grad des Wesens wird in Klammern hinter dem Namen angezeigt
% \tabularnewline typ & \bestTyp & Der \glqq Typ\grqq{} gibt an, zu
% welcher Kategorie die Kreatur gehört, etwa \glqq Dämon\grqq oder
% \glqq Humanoid\grqq{}. Jeder Werte außer \textit{\empty} wird vor dem
% Grad des Wesens innerhalb der gleichen Parenthese abgedruckt.
% \tabularnewline modus & \bestModus & Der \glqq Modus\grqq{} gibt die
% Erscheinungsform des Wesens an. Einige Wesen besitzen verschiedene
% Zustände, etwa \textbf{Werwesen} oder \textbf{Abbadonspinnen}. Bei
% währen dann \textit{in Menschengestalt} bzw.~textit{in Wergestalt}
% sowie \textit{körperlich} bzw.~\textit{auf Seelenreise} angebrachte
% Einträge für |modus|. Jede Angabe außer \textit{\empty} wird
% unmittelbar vor der Klammer in der Titelzeile abgedruckt.
% \tabularnewline variant & \bestVariant & Einige Wesen gibt es in
% verschiedenen Varianten, wobei die Variante jedoch weder (wie bei
% |modus|) eine temporäre Form noch (wie bei |addendum|) ein
% Namenszusatz ist. Beispiele hierfür wären die Angaben der Tierzahlen
% für verschieden große Schwärme oder, etwa bei Krokodilen,
% verschiedene Körpergrößen. Die \glqq Variante\grqq{} wird noch
% zwischen dem |addendum| und dem |modus| in Normalschrift ausgegeben.
% \tabularnewline size & \bestSize & Die Größenordnung des Wesens,
% etwa \textit{Groß} oder \textit{Riesig} (als Regelausdruck). Jede
% Angabe außer \textit{\empty} wird mit einem Gedankenstrich abgetrennt
% hinter dem Grad innerhalb der selben Parenthese abgedruckt.
% \tabularnewline LP & \bestLP & Die Lebenspunkte eines Wesens. Neben
% Würfelausdrücken können hier \textinfty{} (|\textinfty|), * (|*|)
% oder \textstar{} (|\textstar|) sowie - (|-|) geeignete Ausdrücke
% sein. Jede Angabe wird wörtlich übernommen.  \tabularnewline AP &
% \bestAP & Die Lebenspunkte eines Wesens. Neben Würfelausdrücken
% können hier \textinfty{} (|\textinfty|), * (|*|) oder \textstar{}
% (|\textstar|) sowie - (|-|) geeignete Ausdrücke sein. Jede Angabe
% wird wörtlich übernommen.  \tabularnewline RK & \bestRK & Die
% Rüstungsklasse des Wesens, typischerweise eine aus \textbf{OR},
% \textbf{TR}, \textbf{LR}, \textbf{KR}, \textbf{PR}, \textbf{VR} oder
% \textbf{RR}, alternativ auch \textbf{*R}. Jede angabe wird wörtlich
% übernommen.  \tabularnewline EP & \bestEP & Der EP-Multiplikator des
% Wesens wird wörtlich übernommen, aber nur im |\bestiarium|-Layout
% abgedruckt.  \tabularnewline In & \bestIn & Die Intelligenz des
% Wesens wird wörtlich übernommen, aber nur im |\bestiarium|-Layout
% abgedruckt. Angemessen sind Angaben wie \textit{t35} oder
% \textit{m05}.  \tabularnewline St & \bestSt & Die Stärke des Wesens
% wird wörtlich übernommen. Neben Zahlenangaben wie \textit{01} oder
% \textit{100} ist auch ein Stern (* oder \textstar) möglich, etwa für
% Schwarmwesen.  \tabularnewline Gw & \bestGw & Die Stärke des Wesens
% wird wörtlich übernommen.  \tabularnewline B & \bestB & Die
% Bewegungsweite des Wesens wird wörtlich übernommen. Unterschiedliche
% Bewegungsweiten für verschiedene Fortbewegungsarten können mit einem
% einfachen Querstrich (|/|) abgetrennt werden.  \tabularnewline MW &
% \bestMW & Für den Moralwert des Wesens wird jede Angabe außer
% \textit{\empty} als EW interpretiert und in der Form \glqq
% MW+\meta{MW}\grqq{} abgedruckt.  \tabularnewline abwehr &
% \bestAbwehr & Für die Abwehr des Wesens wird jede Angabe als EW
% interpretiert und in der Form \glqq Abwehr+\meta{abwehr}\grqq{}
% abgedruckt. Falls die Angabe keinen Sinn macht, ist ein Stern (*
% oder \textstar) zu verwenden.  \tabularnewline resistenz &
% \bestResistenz & Für die Resistenzen des Wesens wird jede Angabe als
% EW interpretiert und in der Form \glqq
% Resistenz+\meta{resistenz}\grqq{} abgedruckt. Die unterschiedlichen
% Resistenzen sollten in der Form |resistenz=10/12/10| angegeben
% werden. Falls die Angabe keinen Sinn macht, ist ein Stern (* oder
% \textstar) zu verwenden \tabularnewline angriff & \bestAngriff & Die
% Angriffsmöglichkeiten des Wesens sollten hier im für das Regelwerk
% üblichen Format angegeben werden. Durch Komma abgetrennte Einträge
% \textbf{müssen} von geschweiften Klammern|{|\dots|}| umschlossen
% werden. Abwehrwaffen, besondere Angriffsmöglichkeiten und
% \textit{Raufen} sind jeweils mit einem Gedankenstrich \glqq--\grqq{}
% abgetrennt mit einzubinden. Der Ausdruck wird nach dem Wort
% \textbf{Angriff:} in einem eigenen Absatz dargestellt. Nachfolgende
% Zeilen erhalten einen hängenden Einzug.  \tabularnewline zauber &
% \bestZauber & Die Zauberfähigkeiten des Wesens. Jede Angabe außer
% \textit{\empty} wird kursiv hinter dem Erfolgswert dargestellt (siehe
% nächster Eintrag). Wie schon bei |angriff| sind die einzelnen
% Einträge durch Kommas abzutrennen und durch geschweifte Klammern
% |{|\dots|}| zu umschließen.  \tabularnewline zauberEW &
% \bestZauberEW & Der Zaubererfolgswert wird in der Form \glqq
% Zaubern+\meta{zauberEW}:\grqq{} dargestellt und vor den
% Zauberfähigkeiten abgedruckt.  \tabularnewline zauberMore &
% \bestZauberMore & Falls ein Wesen verschiedene Zaubererfolgswerte
% für verschiedene Zauberfähigkeiten besitzt, bietet dieses Feld die
% Möglichkeit, einen weiteren Eintrag zu bilden.  \tabularnewline
% zauberMoreEW & \bestZauberMoreEW & Der Zaubererfolgswert für die
% weiteren Zauberfähigkeiten.  \tabularnewline zauberMMore &
% \bestZauberMMore & Dieses Feld bietet die Möglichkeit für einen
% dritten Eintrag von Zauberfähigkeiten.  \tabularnewline
% zauberMMoreEW & \bestZauberMMoreEW & Der Erfolgswert für den dritten
% \textit{Zaubern}-Eintrag.  \tabularnewline zauberArtig &
% \bestZauberArtig & Angeborene Fähigkeiten, die wie Zauber wirken,
% werden ebenfalls in \textit{Kursiv} dargestellt. Einschränkungen und
% Kommentare sind hier geklammert im Normaldruck üblich, was in der
% Form |zauberArtig={Versetzen \textup{(sich selbst bis zu 50\,km)}}|
% realisiert werden kann.  \tabularnewline zauberArtigEW &
% \bestZauberArtigEW & Die Erfolgswerte für die Anwendung
% zauberartiger, angeborener magischer Fähigkeiten werden in
% (verkehrten) \guillemotright französischen
% Anführungszeichen\guillemotleft{} abgedruckt.  \tabularnewline
% zauberArtigMore & \bestZauberArtigMore & Ein weiteres Feld für mehr
% zauberartige, angeborene Fertigkeiten. Es gelten im Übrigen die
% Anmerkungen zu |zauberArtig|.  \tabularnewline zauberArtigMoreEW &
% \bestZauberArtigMoreEW & Der zugehörige Erfolgswert. Es gelten die
% Anmerkungen zu |zauberArtigEW|.  \tabularnewline zauberArtigMMore &
% \bestZauberArtigMMore & Ein weiteres Feld für mehr zauberartige,
% angeborene Fertigkeiten. Es gelten im Übrigen die Anmerkungen zu
% |zauberArtig|. \tabularnewline zauberArtigMMoreEW &
% \bestZauberArtigMMoreEW & Der zugehörige Erfolgswert. Es gelten die
% Anmerkungen zu |zauberArtigEW|.  \tabularnewline bes & \bestBes &
% Die Besonderheiten werden nach dem Kürzel \glqq\textbf{Bes.:}\grqq{}
% in einem eigenen Absatz dargestellt. Jede Angabe außer \textit{\empty}
% wird wörtlich abgedruckt. Nachfolgende Zeilen erhalten einen
% hängenden Einzug.  \tabularnewline aura & \bestAura & Die Aura des
% Wesens. Jede Angabe außer \textit{\empty} wird nach dem Wort
% \glqq\textsc{Aura}\grqq{} mit einem Gedankenstrich \glqq--\grqq{} an
% die \textbf{Besonderheiten} angehängt.  \tabularnewline vorkommen &
% \bestVork & Das Vorkommen einer Art wird nur im |\bestiarium|-Layout
% dargestellt. Angaben sollten in Stichwortform erfolgen, durch Kommas
% abgetrennt und von geschweiften Klammern |{|\dots|}|
% umschlossen.\tabularnewline energy & \bestEnergie & Falls das
% \textsc{Energie}-Kampfsystem (Hausregelwerk) verwendet wird, kann
% hier die Energie der Lebewesen angegeben werden.
% \tabularnewline \end{longtable} } \end{widecenter}
% \StopEventually{\PrintIndex} \clearpage \section{Implementierung} \begin{macrocode}
\RequirePackage{etex}
\RequirePackage{setspace}
\RequirePackage{xkeyval}
\RequirePackage{etoolbox}
\RequirePackage{graphicx}
\RequirePackage{numprint}
\npthousandsep{\,}
\RequirePackage[framemethod=tikz]{mdframed}
\usetikzlibrary{shadows}
\RequirePackage{longtable}
\RequirePackage{calc}
\RequirePackage{caption}
\RequirePackage{anyfontsize}
\RequirePackage{forloop}
\RequirePackage{etextools}
\RequirePackage{enumitem}
%    \end{macrocode}
%\begin{macro}{\phantomsection}
% Um auch ohne das |hyperref|-Paket fehlerlos Dokumente verarbeiten zu können,
% stellen wir und leeres Kommado von |\phantomsection| zur Verfügung.
%    \begin{macrocode}
\providecommand\phantomsection{}
%    \end{macrocode}
%\end{macro}
%\begin{macro}{\@doendeq}
% Dieses Helferkommando assistiert dabei, die Einrückung nach bestimmten
% Umgebungen zu entfernen.
%    \begin{macrocode}
\newcommand*{\@doendeq}{%
  \everypar{{\setbox\z@\lastbox}\everypar{}}%
}
%    \end{macrocode}
%\end{macro}
%   Die folgenden Counter werden benötigt um die Grund-, und
%   Ausnahmekosten von Zaubern und Fähigkeiten nur aus der Angabe von
%   Standardkosten zu berechnen.
%    \begin{macrocode}
\newcounter{midgard@GrundKosten}
\newcounter{midgard@StandardKosten}
\newcounter{midgard@AusnahmeKosten}
%    \end{macrocode}
%\begin{macro}{\EnergieSystem} Das |midgard|-Paket ist sehr
% flexibel. Das Layout kann an das verwendete Regelsystem angepasst
% werden. Mittels entsprechender Makros kann dann zwischen den
% verwendeten Systemen gewechselt werden, hier im Beispiel zwischen
% dem \textsc{DF4} und dem \textsc{EnergieSystem}.
%    \begin{macrocode}
\newtoggle{useEnergy}
\newtoggle{useInitiative}
\newcommand{\EnergieSystem}{%
  \typeout{package MIDGARD info: Verwende jetzt das Energie-System}%
  \global\toggletrue{useEnergy}%
  \global\toggletrue{useInitiative}%
}
%    \end{macrocode}
%\end{macro}
%\begin{macro}{\DFR} 
% Das Kommando für die Verwendung des Standard-Layouts ist |\DFR|.
%    \begin{macrocode}
\newcommand{\DFR}{%
  \typeout{package MIDGARD info: Verwende jetzt das DFR}%
  \global\togglefalse{useEnergy}%
  \global\togglefalse{useInitiative}%
}
\DFR
%    \end{macrocode}
%\end{macro}
%\begin{macro}{\EnergieOrDFR} 
% Mit diesem Kommando können kurze Textpassagen zwischen zwei Regelversionen
% ein und desselben Dokumentes ausgetauscht werden. Für längere Textpassagen
% sollte das |versions|-Paket zurate gezogen werden.
%    \begin{macrocode}
\newcommand{\EnergieOrDFR}[2]{\iftoggle{useEnergy}{#1}{#2}}
%    \end{macrocode}
%\end{macro}
%
%\begin{macro}{\iquo}
% Mit diesem Helfermakro kann der Ganzzahlquotient zweier Zahlen bestimmt werden.
%    \begin{macrocode}
\newcommand{\iquo}[2]{% integer quotient of #1 divided by #2
  \romannumeral-`\.% for 2-step expansion
  \ifnum#1=0
    \expandafter 0%
  \else
    \the\numexpr(#1-\numexpr(#2-\ifodd \numexpr#2\relax 1 \else 0\fi)/2\relax)/(#2)\expandafter\relax
  \fi
}
%    \end{macrocode}
%\end{macro}
%\begin{macro}{\irem}
% Entsprechend wird der Divisionsrest zweier Zahlen berechnet.
%    \begin{macrocode}
\newcommand{\irem}[2]{% integer remainder of #1 divided by #2
  \the\numexpr #1-#2*\iquo{#1}{#2}\relax}
%    \end{macrocode}
%\end{macro}

%\begin{macro}{\hashstring}
%    Der Befehl |\hashstring| wird intern verwendet, um Querverweise 
%    \textit{automagisch} zu handhaben. Er implementiert ein simples 
%    hashing-verfahren, bei dem jedem Argument zum Zeitpunkt des ersten 
%    Aufrufs eine Nummer zugewiesen wird, zu der das Argument dann zu 
%    jedem späteren Zeitpunkt expandiert.
%    \begin{macrocode} 

\newcounter{@hash@id}
\newcommand\hashstring[1]{%
  \ifcsname\detokenize{hash@#1}\endcsname%
  \else%
  \typeout{info: assigning 'hash:\the@hash@id'  to hash@#1}%
  \expandafter\xdef\csname\detokenize{hash@#1}\endcsname{\the@hash@id}%
  \stepcounter{@hash@id}%
  \fi%
}%
\newcommand\gethash[1]{hash:\csname\detokenize{hash@#1}\endcsname}%

%    \end{macrocode}
%    Um auf das hashing tatsächlich zu verwenden, muss später nurnoch 
%    |hyperref| mit dem expandierten Argument aufgerufen werden.
%    \begin{macrocode}
\newcommand\midgard@htarget[1]{%
  \@ifundefined{hyper@@anchor}{}{%
    \hashstring{midgard@#1}%
    \Hy@raisedlink{\hypertarget{midgard:\gethash{midgard@#1}}{}}%
}}%
\newcommand\midgard@hlink[2]{%
  \@ifundefined{hyper@@anchor}{#1}{%
    \hashstring{midgard@#1}%
    \hyperlink{midgard:\gethash{midgard@#1}}{#2}%
}}%
%    \end{macrocode}
%\end{macro}
%
%\begin{macro}{\midgardskip}
%   Der Befehl |\midgardskip| wird verwendet, um intern an
%   verschiedenen Stellen vertikalen Platz zu erzeugen.
%    \begin{macrocode}
\newlength\midgardskip
\setlength\midgardskip{0.2em}
\newcommand\nutzercode[2]{\expandafter\def\csname midgard@#1@customcode\endcsname{#2}}
%    \end{macrocode}
%\end{macro}
%
%\begin{macro}{spacelesstabbing}
%   Die Umgebung |spacelesstabbing| brauchen wir, da wir die
%   \LaTeX-Umgebung |tabbing| benutzen wollen, um Zauber und
%   Kreaturendaten zu setzen, aber vermeiden wollen, dass die
%   Umgebungen jeweils zusätzliche Abstände zwischen Text und
%   datenblock verursachen.
%    \begin{macrocode}
\newenvironment{spacelesstabbing}{%
\setlength{\topsep}{-\parskip}
\setlength{\parskip}{0pt}
\setlength{\parsep}{0pt}
\setlength{\topskip}{0pt}
\setlength{\partopsep}{0pt}
\setlength{\textfloatsep}{0pt}
  \tabbing%
}
{\endtabbing}

%    \end{macrocode}
% \end{macro}

%\begin{macro}{multicolslongtable}
% In Rollenspielpublikationen sind sehr lange Tabellen nicht
% unüblich. Diese sind oft schwierig zu setzen, da sie ausgesprochen
% schmal sind. Ein mehrspaltiges Tabellenlayout im Stil von
% |multicols| ist hier vorteilhaft, mit normalen \LaTeX-Mitteln aber
% meist nur schwierig zu realisieren. Die
% |multicolslongtable|-Umgebung funktioniert wie |longtable| aus dem
% gleichnamigen Paket, ist aber für den Einsatz innerhalb einer
% |multicols|-Umgebung gedacht.
%    \begin{macrocode}
\newsavebox\ltmcbox
\newenvironment{multicolslongtable}[1]{
  \setbox\ltmcbox\vbox\bgroup
  \col@number\@ne
  \begin{longtable}{#1}
}{
  \end{longtable}
  \unskip
  \unpenalty
  \unpenalty\egroup
  \unvbox\ltmcbox
}  
%    \end{macrocode}
% \end{macro}
%
%\begin{macro}{figurehere}
%    \begin{macrocode}
\newenvironment{figurehere}
               {\par\medskip\noindent\minipage{\linewidth}}
               {\endminipage\par\medskip}
%    \end{macrocode}
%\end{macro}
               
%\begin{macro}{tablehere}
%    \begin{macrocode}
\newenvironment{tablehere}
  {\par\medskip\noindent\minipage{\linewidth}}
  {\endminipage\par\medskip}
%    \end{macrocode}
%\end{macro}

%\begin{macro}{\optendash}
% Ein |\optendash| ist ein |\endash|, der aber als \textit{leader}
% definiert wurde, am Anfang und Ende von Zeilen also unterdrückt
% wird.
%    \begin{macrocode}
\newsavebox\zz
\savebox\zz{--}
\def\optendash{\cleaders\copy\zz\hskip1.9\wd\zz}
%    \end{macrocode}
% \end{macro}
%

%\begin{macro}{unraggedright}
% Mit diesem Kommando kann das Absatzlayout für linksbündige Passagen
% mit kleinen erlaubten Abweichungen von dem rechten Zeilenrand
% eingestellt werden.
%    \begin{macrocode}

\newcommand{\unraggedright}{\raggedright\rightskip=0pt plus 0.2\linewidth\sloppy}
%    \end{macrocode}
% \end{macro}
% \begin{macro}{\text...}
%   Einige mathematische Symbole wie etwa das Unendlich-Zeichen
%   \glqq\textinfty\grqq{} müssen in den Daten von Zaubern und
%   Kreaturen häufig im Text-Modus gesetzt werden. Daher definieren
%   wir hier entsprechende Kommandos. Einige dieser Kommandos werden
%   abhängig vom gewählten Encoding auch von anderen Paketen definiert.
%   Damit sie dennoch wie gewohnt funktionieren, überschreiben wir 
%   bereits erfolge Definitionen am Dokumentanfang.
%    \begin{macrocode}
\AtBeginDocument{
  \global\def\textinfty{\ensuremath{\infty}}
  \global\def\textstar{\ensuremath{\ast}}
  \global\def\texttimes{\ensuremath{\times}}
  \global\def\textdagger{\ensuremath{\dagger}}
  \global\def\textddagger{\ensuremath{\ddagger}}
  \global\def\textsubscript#1{\ensuremath{{}_{\textup{#1}}}}
  \global\def\textsuperscript#1{\ensuremath{{}^{\textup{#1}}}}
  \newcommand{\texthigh}[1]{\ensuremath{{}^{\textup{#1}}}}
  \newcommand{\textlow} [1]{\ensuremath{{}_{\textup{#1}}}}
  \newcommand{\textfrac}[2]{\ensuremath{\frac{#1}{#2}}}
}
%    \end{macrocode}
% \end{macro}
%\begin{macro}{datablist}
% Stellenweise wird das Paket |datablist| verwendet, um 
% Daten von Kreaturen, Zaubern, Fertigkeiten oder Gegenständen
% aus dem Dokument exportieren zu können.
% Damit Dokumente auch prozessiert werden können, ohne das 
% |datablist|-Paket zu verwenden, definieren wir an dieser
% Stelle eine entsprechende Abfrage.
%    \begin{macrocode}
\AtBeginDocument{
  \@ifpackageloaded{datablist}{%
    \providecommand{\ifdatablist}[1]{#1}%
  }{
    \providecommand{\ifdatablist}[1]{}%
  }
}
%    \end{macrocode}
% \end{macro}
% \begin{macro}{\DeclareUnicodeCharacter}
% Bestimmte Sonderzeichen werden oft in Midgard-Dokumenten
% verwendet. Um der Bequemlichkeit willen wollen wir auch das Einfügen
% dieser Sonderzeichen direkt in den \LaTeX Quellcode erlauben.
%    \begin{macrocode}
\@ifpackageloaded{inputenc}{%
\DeclareUnicodeCharacter{00B0}{\ensuremath{^\circ}}
\DeclareUnicodeCharacter{00B2}{\ensuremath{^2}}
\DeclareUnicodeCharacter{00B3}{\ensuremath{^3}}
\DeclareUnicodeCharacter{00B9}{\ensuremath{^1}}
\DeclareUnicodeCharacter{1D9C}{\ensuremath{^\textrm{c}}}
\DeclareUnicodeCharacter{1DA0}{\ensuremath{^\textrm{f}}}
\DeclareUnicodeCharacter{1D56}{\ensuremath{^\textrm{p}}}
\DeclareUnicodeCharacter{02E2}{\ensuremath{^\textrm{s}}}
\DeclareUnicodeCharacter{00BC}{\ensuremath{\nicefrac{1}{4}}}
\DeclareUnicodeCharacter{00D7}{\ensuremath{\times}}
\DeclareUnicodeCharacter{201E}{\glqq}
\DeclareUnicodeCharacter{201C}{\grqq}
\DeclareUnicodeCharacter{221E}{\ensuremath{\infty}}
\DeclareUnicodeCharacter{00AD}{\ignorespaces}
}{}
%    \end{macrocode}
% \end{macro}
% \begin{macro}{\hyphenation}
% Das Midgard-Regelwerk verwendet einige sehr ungewöhnliche Begriffe,
% die sich mit den gängigen Silbentrennungsregeln nicht umbrechen
% lassen. Für diese definieren wir daher die Silbentrennungsregeln
% manuell.
%    \begin{macrocode}
\hyphenation{%
  % general
  schwarz-ma-gisch%
  Dverg-ar Hu-ne Hu-nen Fro-de-volk Spae-folk%
  % Dvergar
  Hrim-dverg Fell-dverg Glod-dverg Shu-Hou Vin-dverg Sael-dverg%
  % Hunen
  Hrim-hune Fjall-hune Har-kar-hune Schu-Dschu-Ren Gjorn-hune Mar-hune%
  % Frodevolk
  I-sing Duns-krat Hark-buar Tschun-Hai-Dsu Myrk-nir Fos-se-mand%
  % Dämonen
  In-dru-val Du-ne-brast Trus-can Ha-le-bant Ka-li-gin Me-gant Co-lus-car%
  % Spaefolk
  Kall-bjargi Ler-gris Ben-Scheng Ka-ring Mok-kurb-jar-gi%
  Sval-grim Haug-ska-di El-dring Chang-Li-Wu Nerf-ling Jarn-bani%
  Snaer-ha-mir Fen-grip Hyrr-bau-ti Schen-Dschi-Gun Thrym-bar Kall-rani%
}
%    \end{macrocode}
% \end{macro}
%\begin{macro}{\HUGE}
% In Midgardpublikationen werden oft sehr große Schriftgrößen verwendet. Daher definieren wir ein neues Schriftgrößenkommando, um diesen Extra-Größen gerecht zu werden.
%    \begin{macrocode}
\newcommand{\HUGE}{\fontsize{36}{36}\selectfont}
%    \end{macrocode}
%\end{macro}
% \begin{macro}{\midgardabenteuer}
%   Die offiziellen Publikationen zeigen auf der dem
%   Inhaltsverzeichnis vorgeschalteten Titelseite häufig einen
%   Schriftzug wie den in Abbildung~\ref{fig:midgardabenteuer}
%   dargestellten. Diesen wollen wir hier nachempfinden.
%    \begin{macrocode}
\newcommand\midgardabenteuer{%
\textbf{\fontsize{36}{0}\selectfont M}
\hspace{-6pt}
\begin{minipage}{65pt}
\vspace{-20pt}
\centering \noindent{\LARGE \textbf{IDGAR}}\\
ABENTEUER
\end{minipage}
\hspace{-6pt}
\textbf{\fontsize{36}{0}\selectfont D}
} 
%    \end{macrocode}
% \end{macro}
% \begin{macro}{\midgardquellenbuch}
%   Ähnlich verhält es sich mit Quellenbüchern.
%    \begin{macrocode}
\newcommand\midgardquellenbuch{%
\textbf{\fontsize{48}{0}\selectfont M}
\hspace{-6pt}
\begin{minipage}[t]{106pt}
\centering \noindent
\vspace{-34pt}
\bfseries
\fontsize{30}{0}\selectfont
\textbf{IDGAR}
\vspace{-6pt}
\par
 \LARGE Quellenbuch
\end{minipage}
\hspace{-6pt}
\textbf{\fontsize{48}{0}\selectfont D}
}
%    \end{macrocode}
%\end{macro}
% \begin{macro}{beispiel}
%   In offiziellen Veröffentlichungen werden Beispiele in leicht grau
%   hinterlegte Rundboxen gesetzt. Dies versuchen wir mit dem Paket
%   |mdframed| nachzubilden. Das Ergebnis wird in
%   Abbildung~\ref{fig:beispiel} dargestellt.
%    \begin{macrocode}
\newmdenv[%
  backgroundcolor=lightgray,%
  roundcorner=10pt,%
%%%  tikzsetting={drop shadow={shadow xshift=1.0ex, shadow yshift=-0.5em, fill=black!50, opacity=1, every shadow }}%
]{beispiel}
%    \end{macrocode}
%\end{macro}
% \begin{macro}{rundbox}
% Auch wollen wir eine Rundbox (Infobox) definieren.
%    \begin{macrocode}
\newenvironment{rundbox}[1][]{
  \begin{mdframed}[roundcorner=10pt,innertopmargin=0.5em,innerbottommargin=0.5em]%
    \ifempty{#1}{}{\noindent {#1}\par}%
    \@afterindentfalse%
    \@afterheading%
}{
  \end{mdframed}
}
%    \end{macrocode}
%\end{macro}
% \begin{macro}{textbox}
% Auch wollen wir eine Textbox (Eckbox) definieren.
%    \begin{macrocode}
\newcommand\textboxstyle{}
\newenvironment{textbox}[1][]{
  \begin{mdframed}[innertopmargin=0.5em,innerbottommargin=0.5em]
    \textboxstyle\textbf{#1}\par
    \@afterindentfalse%
    \@afterheading%
}{
  \end{mdframed}
}
%    \end{macrocode}
% \end{macro}
%
%
% \begin{macro}{\setzeGliederungsEbene}
% Alle für Datenblöcke definiertern Umgebungen fügen dem
% Inhaltsverzeichnis automatisch Einträge hinzu. Die Gliederungsebene,
% auf welcher diese erscheinen, können mit dem Makro
% |\setzeGliederungsEbene{|\meta{Typ}|}{|\meta{Ebene}|}| kontrolliert
% werden.
%    \begin{macrocode}
\def\setzeGliederungsEbene#1#2{%
  \expandafter\def\csname midgard@#1@tocLevel\endcsname{#2}
}
%    \end{macrocode}
% \end{macro}

% \subsection{Das Fantasy Rollenspiel}
%
% \begin{macro}{\geld}
%   Die auf Midgard gängige Währung sind Goldstücke (GS), Silberstücke
%   (SS) und Kupferstücke (KS). Im Standardregelwerk sind die
%   Umrechnungsfaktoren zwischen den einzelnen Münzsorten jeweils 10,
%   d.h.~ein Golstück ist soviel wert wie 10\,SS und ein Silberstück
%   so viel wie 10\,KS. Es gibt jedoch viele Gruppen, die als
%   Hausregeln mit anderen Umrechnungsfaktoren spielen, beispielsweise
%   mit 1:100 (statt 1:10). Um dieser Besonderheit Rechnung zu tragen,
%   definieren wir ein Makro |\geld|, welches eine Anzahl an
%   Kupferstücken entgegennimmt und die jeweils passenden Beträge an
%   Gold- und Silberstücken abdruckt.
%     \begin{macrocode}
\numgdef{\midgard@CopperPerSilver}{10}
\numgdef{\midgard@SilverPerGold  }{10}
\newcommand{\setzeKupferProSilber}[1]{\numgdef{\midgard@CopperPerSilver}{#1}}
\newcommand{\setzeSilberProGold}[1]{\numgdef{\midgard@SilverPerGold}{#1}}
\newcommand{\kupferAlias}{KS}
\newcommand{\silberAlias}{SS}
\newcommand{\goldAlias}  {GS}
\newcommand{\geld}[1]{% compute GP/SP/CP
  \printgold{\iquo{#1}{\midgard@CopperPerSilver*\midgard@SilverPerGold}}% Gold pieces
            {\iquo{\irem{#1}{\midgard@CopperPerSilver*\midgard@SilverPerGold}}{\midgard@CopperPerSilver}}% Silver pieces
            {\irem{\irem{#1}{\midgard@CopperPerSilver*\midgard@SilverPerGold}}{\midgard@CopperPerSilver}}% Copper pieces
}
\newcommand{\printgold}[3]{% Prints GP/SP/CP
  % #1 = number of gold pieces
  % #2 = number of silver pieces
  % #3 = number of copper pieces
  \ifnum #1 > 0 
    \numprint{#1}\,\goldAlias
  \fi
  \ifnum \numexpr #1*#2 > 0 
    \thinspace%
  \fi
  \ifnum #2 > 0 
    #2\,\silberAlias
  \fi
  \ifnum \numexpr #2*#3 > 0 
    \thinspace%
    \else \ifnum \numexpr #1*#3 > 0 
      \thinspace%
    \fi
  \fi
  \ifnum #3 > 0 
    #3\,\kupferAlias
  \fi
}
%    \end{macrocode}
% \end{macro}
% \begin{macro}{\fertigkeit}
%   Allgemeine Fertigkeiten werden im \textsc{Dfr4} auf ganz
%   charakteristische Weise gesetzt. Zunächst definieren wir unter
%   Verwendung des Paketes |xkeyval| lokale Kommandos, welche (als
%   \glqq\textit{Variablen}\grqq) fungieren und die dem Kommando im
%   optionalen Argument übergebenen Werte speichern.
%    \begin{macrocode}
\setzeGliederungsEbene{fertigkeit}{subsubsection}
\define@key{fertigkeit}{name}{\def\fertigkeitName{#1}}
\define@key{fertigkeit}{entry}{\def\fertigkeitEntry{#1}}
\define@key{fertigkeit}{kategorie}{\def\fertigkeitKategorie{#1}}
\define@key{fertigkeit}{leitAtt}{\def\fertigkeitLeitAtt{#1}}
\define@key{fertigkeit}{anforderungen}{\def\fertigkeitAnforderungen{#1}}
\define@key{fertigkeit}{EW}{\def\fertigkeitEW{#1}}
\define@key{fertigkeit}{ungelernt}{\def\fertigkeitUngelernt{#1}}
\define@key{fertigkeit}{spielbeginn}{\def\fertigkeitSpielbeginn{#1}}
\define@key{fertigkeit}{maximum}{\def\fertigkeitMaximum{#1}}
\define@key{fertigkeit}{FP}{\def\fertigkeitFP{#1}}
\define@key{fertigkeit}{G}{\def\fertigkeitG{#1}}
\define@key{fertigkeit}{S}{\def\fertigkeitS{#1}}
\define@key{fertigkeit}{A}{\def\fertigkeitA{#1}}
\define@key{fertigkeit}{key}{\def\fertigkeitKey{#1}}
\define@key{fertigkeit}{anmerkung}{\def\fertigkeitAnmerkung{#1}}
%    \end{macrocode}
% Wir definieren zusätzlich \textit{default}-Werte für die Variablen,
% für den Fall, dass sie nicht aufgerufen werden.
%    \begin{macrocode}
\presetkeys{fertigkeit}{name=Fähigkeit, entry=\empty, kategorie=\empty, leitAtt=\empty, 
  anforderungen=\empty, EW=\empty, ungelernt=\empty, spielbeginn=\empty, maximum=\empty, FP=hide, 
  G=niemand, S=niemand, A=niemand, key=ability, anmerkung=\relax}{}
%    \end{macrocode}
% Das tatsächliche Kommando beginnt mit der Empfehlung zu einem
% Seitenumbruch, um unschöne Umbrüche im Layout selbst zu vermindern.
%    \begin{macrocode}
\newcommand\midgard@fertigkeit@customcode\relax%
\newcommand*\fertigkeit[2][]{%
\vspace{0.5em plus 0.5em}
\pagebreak[2]\begingroup%
%    \end{macrocode}
% Anschließend wird der Befehl |setkeys| aufgerufen, um auf die zuvor
% definierten Kommandos zurückgreifen zu können.
%    \begin{macrocode}
\setkeys{fertigkeit}{#1}{%
%    \end{macrocode}
% Wenn |\fertigkeitFP| \textit{hide} ist, müssen die FP nicht ausgelesen
% werden. Andernfalls werden Counter für Standard- und Ausnahmekosten
% definiert und nach der für Fähigkeiten üblichen Formel aus den
% Standardkosten berechnet.
%    \begin{macrocode}
\ifdefstring{\fertigkeitFP}{hide}{}{%
  \setcounter{midgard@StandardKosten}{\fertigkeitFP}%
  \setcounter{midgard@GrundKosten}{\value{midgard@StandardKosten}/2}%
  \setcounter{midgard@AusnahmeKosten}{\value{midgard@StandardKosten}*2}%
}%
\phantomsection%
\setlength{\parindent}{0pt}%
\setlength{\parskip}{0pt}%
\ifdefstring{\fertigkeitEntry}{\empty}{%
  \addcontentsline{toc}{\midgard@fertigkeit@tocLevel}{\fertigkeitName}%
}{%
  \addcontentsline{toc}{\midgard@fertigkeit@tocLevel}{\fertigkeitEntry}%
}%
\expandnext{\midgard@htarget}{\fertigkeitName}\textbf{\fertigkeitName}%
\ifboolexpr{ test {\ifdefstring{\fertigkeitFP}{hide}} or test {\ifdefstring{\fertigkeitEW}{\empty}} }{}{+\fertigkeitEW}%
\ifdefstring{\fertigkeitKategorie}{\empty}{}{~(\fertigkeitKategorie)}%
\ifdefstring{\fertigkeitAnmerkung}{\relax}{}{ \optendash{} \textit{\fertigkeitAnmerkung}}%
\ifboolexpr{ test {\ifdefstring{\fertigkeitFP}{hide}} or test {\ifdefstring{\fertigkeitEW}{\empty}} }{%
  \ifdefstring{\fertigkeitLeitAtt}{\empty}{}{\hfill\fertigkeitLeitAtt\hfill}%
  \ifdefstring{\fertigkeitAnforderungen}{\empty}{}{\hfill{\fertigkeitAnforderungen}\hfill}~%
}{%
  \ifdefstring{\fertigkeitUngelernt}{\empty}{}{\hfill ungelernt+(\fertigkeitUngelernt)}%
  \par\vspace{\midgardskip}
  \begin{small}
    \ifdefstring{\fertigkeitLeitAtt}{\empty}{}{\textbf{\fertigkeitLeitAtt}}%
    \ifdefstring{\fertigkeitAnforderungen}{\empty}{}{%
      \ifdefstring{\fertigkeitLeitAtt}{\empty}{}{, }{\fertigkeitAnforderungen}%
    }%
    \hfill~\linebreak[0] \mbox{Erfolgswert+\fertigkeitEW}\ifdefstring{\fertigkeitSpielbeginn}{\empty}{}{\hfill(+\fertigkeitSpielbeginn/+\fertigkeitMaximum)}%
  \end{small}%
}\par%
\ifdefstring{\fertigkeitFP}{hide}{}{%
  \begin{small}%
  \vspace{\midgardskip}
  \ifdefstring{\fertigkeitG}{niemand}%
              {niemand }%
              {\textbf{\themidgard@GrundKosten :} \nolinebreak \fertigkeitG} -- %
              \ifdefstring{\fertigkeitS}{niemand}%
                          {niemand }%
                          {\textbf{\themidgard@StandardKosten :} \nolinebreak \fertigkeitS} -- %
                          \ifdefstring{\fertigkeitA}{niemand}%
                                      {niemand }%
                                      {\textbf{\themidgard@AusnahmeKosten :} \nolinebreak \fertigkeitA}
                                      \par
    \end{small}%
    \vspace{\midgardskip}
}%
\expandafter\ifcsdef{\fertigkeitKey exportvalues}{%
  \ifdatablist{\adddataline{\fertigkeitKey}{\csname \fertigkeitKey exportvalues\endcsname}}%
}{}%
}\endgroup\par%
\midgard@fertigkeit@customcode%
\@afterindentfalse%
\@afterheading%
}
%    \end{macrocode}
% \end{macro}
%
% \begin{macro}{steigertabelle}
%    \begin{macrocode}
\newenvironment{steigertabelle}[1]{%
  \begin{center}
\begin{minipage}{.8\linewidth}
\setlength{\parindent}{0pt}
\setlength{\parskip}{0pt}
  \setkeys{fertigkeit}{#1}%
  \textbf{\fertigkeitName+\fertigkeitEW}\par
  \hspace{1em}\begin{tabular}{c >{\bfseries}r}
}{%
  \end{tabular}\par
  \begin{description}[font=\normalfont\itshape,noitemsep,nolistsep]
    \ifdefstring{\fertigkeitG}{niemand}{}{\item[Grund:]    \fertigkeitG}
    \ifdefstring{\fertigkeitS}{niemand}{}{\item[Standard:]    \fertigkeitS}
    \ifdefstring{\fertigkeitA}{niemand}{}{\item[Ausnahme:]    \fertigkeitA}
  \end{description}
\end{minipage}
\end{center}
}
%    \end{macrocode}
% \end{macro}
% 
% \begin{macro}{\waffenfertigkeit}
%   Waffenfertigkeiten werden im \textsc{Dfr4} auf ganz
%   charakteristische Weise gesetzt. Zunächst definieren wir lokale
%   Kommandos.
%    \begin{macrocode}
\define@key{weapon}{name}{\def\weaponName{#1}}
\define@key{weapon}{entry}{\def\weaponEntry{#1}}
\define@key{weapon}{schaden}{\def\weaponDamage{#1}}
\define@key{weapon}{energie}{\def\weaponEnergy{#1}}
\define@key{weapon}{reichweite}{\def\weaponRange{#1}}
\define@key{weapon}{anfangserfolgswert}{\def\weaponStartEW{#1}}
\define@key{weapon}{voraussetzungen}{\def\weaponPremise{#1}}
\define@key{weapon}{grundkenntnisse}{\def\weaponBasics{#1}}
\define@key{weapon}{beispiele}{\def\weaponExamples{#1}}
\define@key{weapon}{schwierigkeit}{\def\weaponDifficulty{#1}}
\define@key{weapon}{FP}{\def\weaponFP{#1}}
\define@key{weapon}{G}{\def\weaponG{#1}}
\define@key{weapon}{S}{\def\weaponS{#1}}
\define@key{weapon}{A}{\def\weaponA{#1}}
\define@key{weapon}{key}{\def\weaponKey{#1}}
%    \end{macrocode}
% Wir definieren \textit{default}-Werte für die Variablen.
%    \begin{macrocode}
\presetkeys{weapon}{name=Waffe, entry=\empty, energie=\empty, schaden=-, reichweite=, beispiele=\empty,
  voraussetzungen={St\,01, Gs\,01}, grundkenntnisse=-, anfangserfolgswert=4,
  schwierigkeit=normal, FP=hide, G=niemand, S=niemand, A=niemand, key=weaponability}{}
%    \end{macrocode}
% Das tatsächliche Kommando beginnt mit der Empfehlung zu einem
% Seitenumbruch, um unschöne Umbrüche im Layout selbst zu vermindern.
%    \begin{macrocode}
\newcommand*\waffenfertigkeit[2][]{%
\pagebreak[2]
%    \end{macrocode}
% Anschließend wird der Befehl |setkeys| aufgerufen, um auf die zuvor
% definierten Kommandos zurückgreifen zu können.
%    \begin{macrocode}
\setkeys{weapon}{#1}{
%    \end{macrocode}
% Wenn |\weaponFP| \textit{hide} ist, müssen die FP nicht ausgelesen
% werden. Andernfalls werden Counter für Standard- und Ausnahmekosten
% definiert und nach der für Fähigkeiten üblichen Formel aus den
% Standardkosten berechnet.
%    \begin{macrocode}
\ifdefstring{\weaponFP}{hide}{}{
  \setcounter{midgard@StandardKosten}{\weaponFP}%
  \setcounter{midgard@GrundKosten}{\value{midgard@StandardKosten}/2}%
  \setcounter{midgard@AusnahmeKosten}{\value{midgard@StandardKosten}*2}%
}%
\phantomsection%
\ifdefstring{\weaponEntry}{\empty}{%
  \addcontentsline{toc}{\midgard@fertigkeit@tocLevel}{\weaponName}
}{%
  \addcontentsline{toc}{\midgard@fertigkeit@tocLevel}{\weaponEntry}
}
\par\vspace{0.2em plus 0.3em minus 0.1em}%
\raggedright
\par
\noindent\textbf{\weaponName}~\mbox{(\weaponDamage{}%
\ifboolexpr{ togl {useEnergy} and not test {\ifdefstring{\weaponEnergy}{\empty}}}{,~\textsc{E}\,\weaponEnergy}{}%
)}\hfil\mbox{}\nobreak\hfill\nobreak\mbox{\weaponRange{}}\par% 
\begin{small}
\noindent\weaponPremise \hspace{1cm} Erfolgswert+\weaponStartEW\par%
\noindent\textbf{\weaponBasics} \ifdefstring{\weaponExamples}{\empty}{}{\textit{(\weaponExamples)}}
\ifdefstring{\weaponFP}{hide}{%
}{%
  \ifdefstring{\weaponDifficulty}{hide}{}{\newline\noindent\textit{Schwierigkeit:} \weaponDifficulty}%
\\%
  \noindent%
  \ifdefstring{\weaponG}{niemand}%
  {niemand }%
  {\textbf{\themidgard@GrundKosten :} \nolinebreak \weaponG} -- %
  \ifdefstring{\weaponS}{niemand}%
  {niemand }%
  {\textbf{\themidgard@StandardKosten :} \nolinebreak \weaponS} -- %
  \ifdefstring{\weaponA}{niemand}%
  {niemand }%
  {\textbf{\themidgard@AusnahmeKosten :} \nolinebreak \weaponA}\par
}
\end{small}
\expandafter\ifcsdef{\weaponKey exportvalues}{%
  \ifdatablist{\adddataline{\weaponKey}{\csname \weaponKey exportvalues\endcsname}}
}{}%
}
\vspace{0.2em plus 0.3em minus 0.1em}
\@afterindentfalse%
\@afterheading%
}
%    \end{macrocode}
% \end{macro}
% \begin{macro}{\weapon}
%   Aus Gründen der Kompatiblität und Bequemlichkeit definieren wir
%   das Kommand |\weapon| als identisch mit |\waffenfertigkeit|.
%    \begin{macrocode}
\let\weapon\waffenfertigkeit
%    \end{macrocode}
% \end{macro}
%
% \subsection{Arkanum}
%

% \begin{macro}{\img...}  Das \textsc{Midgard}-Regelwerk verwendet
% einige Sonderzeichen, welche als Vektorgraphiken mit freundlicher
% Genehmingung des Verlages mit diesem Paket gemeinsam ausgeliefert
% werden. Diese können mit den hier aufgeführten Kommandos im Text
% eingebunden werden.  \begin{macrocode}
\newcommand\imgMagisch {\includegraphics[height=0.7em]{pics/scroll}}
\newcommand\imgModus   {\includegraphics[height=0.7em]{pics/box}}
\newcommand\imgObjectum{\includegraphics[height=0.7em]{pics/arrow}}
\newcommand\imgFeuer   {\includegraphics[height=0.7em]{pics/fire}}
\newcommand\imgWasser  {\includegraphics[height=0.7em]{pics/water}}
\newcommand\imgErde    {\includegraphics[height=0.7em]{pics/earth}}
\newcommand\imgLuft    {\includegraphics[height=0.55em]{pics/air}}
\newcommand\imgEis     {\includegraphics[height=0.7em]{pics/ice}}
\newcommand\imgHolz    {\includegraphics[height=0.7em]{pics/wood}}
\newcommand\imgMetall  {\includegraphics[height=0.7em]{pics/metal}}
\newcommand\imgMagan   {\includegraphics[height=0.7em]{pics/magic}}
%    \end{macrocode}
% \end{macro}
% \begin{macro}{\zaubervariante}
%   Das Kommando |\zaubervariante| soll benutzt werden um am Ende
%   einer Zauberbeschreibung die Variationen des Zaubers für
%   verschiedene Zaubererklassen genauer beschreiben, etwa |\zaubervariante{Thaumaturgie}| für einen Abschnitt, der mit \glqq\textbf{\textsc{Thaumaturgie}:} \dots\grqq{} beginnt. Zur Implementierung
%   vermeiden wir einen Einzug und setzen das Argument mit |\textbf|
%   \textbf{fett} und mit |\textsc| in \textsc{Kapitälchen}.
%    \begin{macrocode}
\newcommand{\zaubervariante}[1]{\noindent\textbf{\textsc{#1}: }}
%    \end{macrocode}
% \end{macro}
% \begin{macro}{thaumagralzauber}
% Eine besondere Form der Zaubervariante ist der Thaumagralzauber. Dieser erhält einen eigenen Datenblock.
%    \begin{macrocode}
\NewDocumentEnvironment{thaumagralzauber}{O{0} O{Thaumagral}}{%
  \par\pagebreak[2]\par
  \noindent\textbf{#2} ABW+#1%
  \setlength\parindent{1em}%
  \setlength\parskip{1pt}%
  \nopagebreak\par\nopagebreak%
  \renewcommand{\item}[1][]{\par\normalfont ##1\hfill\bfseries}%
}{
  \unskip\unskip
  % these two unskips are needed to remove trailing spaces/newlines
  % see http://tex.stackexchange.com/questions/166158/why-is-unskip-not-working-here
  \ignorespacesafterend\par\noindent\aftergroup\@doendeq%
}
%    \end{macrocode}
% \end{macro}
%
% \begin{macro}{\zauberref}
%   Mit dem |\zauberref|-Kommando erlauben wir einfaches referenzieren auf Zauber.
%   \begin{macrocode}
\newcommand*\zauberref[2][]{{\itshape\ifstrempty{#1}{\protect\midgard@hlink{#2}{#2}}{\protect\midgard@hlink{#1}{#2}}}}%%
\newcommand*\fertigkeitsref[2][]{{\itshape\ifstrempty{#1}{\protect\midgard@hlink{#2}{#2}}{\protect\midgard@hlink{#1}{#2}}}}%%
%   \end{macrocode}
% \end{macro}
% \begin{macro}{\zauber}
%   Um die Umgebung für das Arkanums-Artige Layout neuer Zauber zu
%   setzen, wird das Paket |xkeyval| verwendet. Zunächst werden
%   innerhalb des |\zauber|-Kommandos lokale Kommandos definiert, um
%   auf die Werte der einzlenen Felder zuzugreifen.
%    \begin{macrocode}
\setzeGliederungsEbene{zauber}{subsubsection}
\define@key{zauber}{name}{\def\zaubername{#1}}
\define@key{zauber}{entry}{\def\zauberentry{#1}}
\define@key{zauber}{anmerkung}{\def\zauberAnmerkung{#1}}
\define@key{zauber}{untertitel}{\def\zauberUntertitel{#1}}
\define@key{zauber}{scroll}{\def\zauberscroll{#1}}
\define@key{zauber}{th}{\def\zauberth{#1}}
\define@key{zauber}{art}{\def\zauberart{#1}}
\define@key{zauber}{artshort}{\def\zauberartshort{#1}}
\define@key{zauber}{kurztext}{\def\zaubershortdescription{#1}}
\define@key{zauber}{stufe}{\def\zauberstufe{#1}}
\define@key{zauber}{material}{\def\zaubermaterial{#1}}
\define@key{zauber}{matpreis}{\def\zaubermatpreis{#1}}
\define@key{zauber}{matannotation}{\def\zaubermatannotation{#1}}
\define@key{zauber}{koeder}{\def\beschwoerungskoeder{#1}}
\define@key{zauber}{koederpreis}{\def\beschwoerungskoederpreis{#1}}
\define@key{zauber}{feld}{\def\beschwoerungsfeld{#1}}
\define@key{zauber}{feldpreis}{\def\beschwoerungsfeldpreis{#1}}
\define@key{zauber}{prozess}{\def\zauberprozess{#1}\def\beschwoerungsecprozess {#1}}
\define@key{zauber}{agens}  {\def\zauberagens  {#1}\def\beschwoerungsecagens   {#1}}
\define@key{zauber}{reagens}{\def\zauberreagens{#1}\def\beschwoerungsecreagens {#1}}
\define@key{zauber}{secprozess} {\def\beschwoerungsecprozess {#1}}
\define@key{zauber}{secagens}   {\def\beschwoerungsecagens   {#1}}
\define@key{zauber}{secreagens} {\def\beschwoerungsecreagens {#1}}
\define@key{zauber}{primprozess}{\def\beschwoerungprimprozess{#1}}
\define@key{zauber}{primagens}  {\def\beschwoerungprimagens  {#1}}
\define@key{zauber}{primreagens}{\def\beschwoerungprimreagens{#1}}
\define@key{zauber}{showformula}{\def\showformula{#1}}
\define@key{zauber}{AP}{\def\zauberAP{#1}}
\define@key{zauber}{GZ}{\def\zauberGZ{#1}}
\define@key{zauber}{Zd}{\def\zauberZd{#1}}
\define@key{zauber}{E}{\def\zauberE{#1}}
\define@key{zauber}{Rw}{\def\zauberRw{#1}}
\define@key{zauber}{Rwshort}{\def\zauberRwshort{#1}}
\define@key{zauber}{Wz}{\def\zauberWz{#1}}
\define@key{zauber}{Wb}{\def\zauberWb{#1}}
\define@key{zauber}{Wd}{\def\zauberWd{#1}}
\define@key{zauber}{Ur}{\def\zauberUr{#1}}
\define@key{zauber}{FP}{\def\zauberFP{#1}}
\define@key{zauber}{G}{\def\zauberG{#1}}
\define@key{zauber}{S}{\def\zauberS{#1}}
\define@key{zauber}{A}{\def\zauberA{#1}}
\define@key{zauber}{spelltype}{\def\zauberKey{#1}}
%    \end{macrocode}
% Wir definieren \textit{default}-Werte für die Variablen.
%    \begin{macrocode}
\presetkeys{zauber}{name=Zauber, entry=\empty, scroll=true, th=\empty, kurztext={},
  anmerkung=\empty, untertitel=\empty, art=Zauber, stufe=1, spelltype=spell,
  material=\empty, matpreis=\empty, matannotation=\empty, 
  koeder=\empty, koederpreis=\empty, feld=\empty, feldpreis=\empty,
  prozess=Prozess, agens=Agens, reagens=Reagens,
  AP=0, GZ=hide, Zd=keine, Rw=-, E=hide, Wz=-, Wb=-, Wd=-, Ur={n.\,A.},
  FP=hide, G=niemand, S=niemand, A=niemand}{}
 %  ,  artshort=\empty, Rwshort=\empty}{}
%    \end{macrocode}
% Das |zauber|-Kommando wird eingeleitet mit etwas vertikalem Platz,
% der nach Bedarf (z.\,B.~am Seitenanfang) nicht dargestellt
% wird. Anschließend wird der Befehl |setkeys| aufgerufen, um auf die
% zuvor definierten Kommandos zurückgreifen zu können. Weiterhin geben
% wir iene Empfehlung zum Seitenumbruch unmittelbar vor dem
% Datenblock. Außerdem werden ein paar charakteristische
% Textsatzlängen definiert, um einen hängenden Einzug von überlangen
% Zeilen zu erzeugen..
%    \begin{macrocode}
\newcommand*\zauber[2][]{%
  \par\pagebreak[2]\addvspace{1em}
\begingroup
\setkeys{zauber}{#1}{
%    \end{macrocode}
% Wenn |\zauberFP| \textit{hide} ist, müssen die FP nicht ausgelesen
% werden. Andernfalls werden Counter für Standard- und Ausnahmekosten
% definiert und nach der für Zauber üblichen Formel aus den
% Standardkosten berechnet.
%    \begin{macrocode}
\ifdefstring{\zauberFP}{hide}{}{
  \setcounter{midgard@StandardKosten}{\zauberFP}%
  \setcounter{midgard@GrundKosten}{\value{midgard@StandardKosten}/2}%
  \setcounter{midgard@AusnahmeKosten}{\value{midgard@StandardKosten}*5}%
}%
%    \end{macrocode}
% Wenn |\zauberentry| \textit{\empty}, wird der |\zaubername| auf der
% Gliederungsebene |subsubsection| zum Inhaltsverzeichnis
% hinzugefügt. Andernfalls geschieht dasselbe mit |\zauberentry|.
%    \begin{macrocode}
\phantomsection%
\ifdefstring{\zauberentry}{\empty}{%
  \addcontentsline{toc}{\midgard@zauber@tocLevel}{\zaubername}
}{%
  \addcontentsline{toc}{\midgard@zauber@tocLevel}{\zauberentry}
}
%    \end{macrocode}
% Zunächst wird der Name des Zaubers gedruckt. Dann wird zunächst
% |\hfill| benutzt, um an das Ende der Zeile vorzurücken. 
%    \begin{macrocode}
\noindent\makebox[\linewidth]{\parbox[b]{0.85\linewidth}{%
    \unraggedright\expandnext{\midgard@htarget}{\zaubername}\noindent \bfseries \large  \zaubername%
  }%
\hfill%
\ifdefstring{\zauberscroll}{true}%
           {}%
           {\imgMagisch}%
\hfill%
\ifdefstring{\zauberth}{\empty}%
           {}%
           {\textbf{\zauberth}}%
}\par%
\setlength{\parindent}{-1ex}%
\setlength{\leftskip}{1ex}%
\setlength{\rightskip}{0pt plus 0.2\linewidth}%
\setlength{\parskip}{0pt}%
\par
%    \end{macrocode}
% Falls |\zauberAnmerkung| gesetzt ist, wird sie jetzt klein, fett und
% kursiv abgedruckt.
%    \begin{macrocode}
\ifdefstring{\zauberAnmerkung}{\empty}{}%
{{\par\footnotesize\bfseries\itshape\zauberAnmerkung\par}}%
%    \end{macrocode}
% Falls |\zauberUntertitel| gesetzt ist, wird sie jetzt klein, fett und
% kursiv abgedruckt.
%    \begin{macrocode}
\ifdefstring{\zauberUntertitel}{\empty}{}%
{{\nopagebreak\par\itshape\zauberUntertitel\par}}%
%    \end{macrocode}
% Wir unterdrücken einen Zeilenumbruch unmittelbar nach der
% Titelzeile. 
%    \begin{macrocode}
\nopagebreak\par%
%    \end{macrocode}
% Die Zauberart wird wörtlich übernommen. Ist die Stufe
% \textit{GM}, so wird der Text \glqq der Großen Magie\grqq{}
% nachgestellt, andernfalls \glqq der Stufe |\zauberstufe|\grqq{}
%    \begin{macrocode}
\vspace{\midgardskip}\par
\zauberart~%
\ifdefstring{\zauberstufe}{GM}%
           {der Großen Magie}%
           {der Stufe \zauberstufe}%
\par
%    \end{macrocode}
% Falls |\zaubermaterial| \textit{\empty} ist, gibt es kein
% Zaubermaterial. Wir unterdrücken etwas vertikalen Platz und machen
% sofort mit der Zauberformel weiter. Andernfalls drucken wir zunächst
% das Material, und dann -- wiederum nur falls diese nicht
% \textit{\empty} sind -- |\zaubermatpreis| und |\zaubermatannotation|,
% ersteres geklammert, letzteres durch Semikolon abgetrennt.
%    \begin{macrocode}
\def\zaubermattotal{}%
\ifdefstring{\zaubermaterial}{\empty}%
            {}%
            {\appto\zaubermattotal{\textit{\zaubermaterial}}%
              \ifdefstring{\zaubermatpreis}{\empty}%
                          {}%
                          {\appto\zaubermattotal{~(\zaubermatpreis)}%
                            \ifdefstring{\zaubermatannotation}{\empty}%
                                        {}%
                                        {\appto\zaubermattotal{; \zaubermatannotation}}}}%
\zaubermattotal\par%
%    \end{macrocode}
% Die Zauberformel wird mit etwas Abstand gedruckt, dazwischen jeweils
% die Symbole \imgModus{} und \imgObjectum.
%    \begin{macrocode}
\vspace{\midgardskip}
\zauberprozess~%
\imgModus~%
\zauberagens~%
\imgObjectum~%
\zauberreagens\par\vspace{\parskip}%
\vspace{\midgardskip}
%    \end{macrocode}
% In einer kompaktifizierten |tabbing|-Umgebung drucken wir die
% wichtigen Zauberdaten jeweils, falls diese nicht auf \textit{hide}
% gesetzt sind.
%    \begin{macrocode}
\begin{spacelesstabbing}%
\textbf{Wirkungsbereich:}\hspace{1ex} \= anytext \kill
\ifdefstring{\zauberAP}{hide}{}{\textbf{AP-Verbrauch:} \> \zauberAP\\}%
\ifdefstring{\zauberGZ}{hide}{}{\textbf{Gefährdungszahl:} \> \zauberGZ\\}%
\ifboolexpr{ test {\ifdefstring{\zauberZd}{hide}} or ( togl {useEnergy} and not test {\ifdefstring{\zauberE}{hide}} ) }{}{\textbf{Zauberdauer:} \> \zauberZd\\}%
\ifboolexpr{ togl {useEnergy} and not test {\ifdefstring{\zauberE}{hide}}}{\textbf{Energie} \> \zauberE\\}{}%
\ifdefstring{\zauberRw}{hide}{}{\textbf{Reichweite:} \> \zauberRw\\}%
\ifdefstring{\zauberWz}{hide}{}{\textbf{Wirkungsziel:} \> \zauberWz\\}%
\ifdefstring{\zauberWb}{hide}{}{\textbf{Wirkungsbereich:} \> \zauberWb\\}%
\ifdefstring{\zauberWd}{hide}{}{\textbf{Wirkungsdauer:} \> \zauberWd\\}%
\ifdefstring{\zauberUr}{hide}{}{\textbf{Ursprung:} \> \zauberUr\\}
\end{spacelesstabbing}\vspace{-1.0em}\par%
%    \end{macrocode}
% Wir unterdrücken wiederum etwas vertikalen Platz und drucken dann,
% jeweils falls gesetzt, die Lernkosten. Hierbei ist wichtig, dass wir
% den Fall \textit{niemand} gesondert behandeln, da in diesem Fall
% auch die Kosten selbst nicht abgedruckt werden sollen.
%    \begin{macrocode}
\vspace{1.0\midgardskip}\par
\ifdefstring{\zauberFP}{hide}{}{
  \ifdefstring{\zauberG}{niemand}%
  {niemand }%
  {\textbf{\themidgard@GrundKosten :} \nolinebreak \zauberG} -- %
  \ifdefstring{\zauberS}{niemand}%
  {niemand }%
  {\textbf{\themidgard@StandardKosten :} \nolinebreak \zauberS} -- %
  \ifdefstring{\zauberA}{niemand}%
  {niemand }%
  {\textbf{\themidgard@AusnahmeKosten :} \nolinebreak \zauberA}\par%
  \vspace{1.6\midgardskip}
}%
\expandafter\ifcsdef{\zauberKey exportvalues}{%
  \ifdatablist{\adddataline{\zauberKey}{\csname \zauberKey exportvalues\endcsname}}%
}{}%
}\endgroup%
%    \end{macrocode}
% Abschließend sorgen wir für einen ausreichenden Abstand und
% verhindern die Einrückung des auf den Datenblock folgenden Absatzes.
%    \begin{macrocode}
\@afterindentfalse%
\@afterheading%
}%
%    \end{macrocode}
%\end{macro}
%\begin{macro}{\zaubersalz}
% Aus Bequemlichkeitsgründen definieren wir auch ein Kommando für Zaubersalze, das intern jedoch das Kommando für Zauber aufruft -- mit den voreingestellten Standard-Spieldaten von Zaubersalzen.
%    \begin{macrocode}
\newcommand*\zaubersalz[2][]{
\zauber[%
    spelltype=salt,
    name=Zaubersalz,%
    art=Gedankenzauber,% 
    stufe=1,%
    AP=1,%
    Rw={500\,m},%
    Zd={1\,sec},%
    Rw={500\,m},%
    material=Zaubersalz,%
    matpreis={\geld{20}},%
    Ur={dämonisch},%
    #1%
]{#2}%
}
%    \end{macrocode}
%\end{macro}
%\begin{macro}{\zauberlied}
% Ebenfalls aus Bequemlichketisgründen definieren wir auch ein
% Kommando für Zauberlieder, das intern jedoch das Kommando für Zauber
% aufruft -- mit den voreingestellten Standard-Spieldaten von
% Zauberliedern.
%    \begin{macrocode}
\newcommand*\zauberlied[2][]{%
\zauber[%
    spelltype=music,
    name=Zauberlied,%
    Zd=hide,%
    Rw=-,%
    AP=2,%
    Wz=Geist,%
    art=Zauberlied,%
    Ur=bardisch,%
    #1%
]{#2}%
}
%    \end{macrocode}
%\end{macro}
%\begin{macro}{\zaubertanz}
% Ebenfalls aus Bequemlichketisgründen definieren wir auch ein Kommando für Zaubertänze, das intern jedoch das Kommando für Zauber aufruft -- mit den voreingestellten Standard-Spieldaten von Zaubertänze.
%    \begin{macrocode}
\newcommand*\zaubertanz[2][]{%
\zauber[%
    spelltype=dance,
    name=Zaubertanz,%
    art=Zaubertanz,
    Rw={0\,m},%
    Ur=bardisch,%
    #1%
]{#2}%
}
%    \end{macrocode}
%\end{macro}
%\begin{macro}{\beschwoerung}
%   Völlig analog definieren wir ein Kommando für das Layout von Beschwörungen.
%    \begin{macrocode}
\newcommand*\beschwoerung[2][]{%
  \pagebreak[2]\addvspace{1em}
\begingroup
\presetkeys{zauber}{name=Beschwörung, scroll=false, art=Beschwörung, spelltype=incantation,%
  primprozess=Verändern, primagens=Magan, primreagens=Metall,%
  secprozess=Beherrschen, secagens=Magan, secreagens=Wasser,%
  showformula=\empty%
}{}
\setkeys{zauber}{#1}{
%    \end{macrocode}
% Wenn |\zauberFP| \textit{hide} ist, müssen die FP nicht ausgelesen
% werden. Andernfalls werden Counter für Standard- und Ausnahmekosten
% definiert und nach der für Zauber üblichen Formel aus den
% Standardkosten berechnet.
%    \begin{macrocode}
\ifdefstring{\zauberFP}{hide}{}{
  \setcounter{midgard@StandardKosten}{\zauberFP}%
  \setcounter{midgard@GrundKosten}{\value{midgard@StandardKosten}/2}%
  \setcounter{midgard@AusnahmeKosten}{\value{midgard@StandardKosten}*5}%
}%
%    \end{macrocode}
% Wenn |\zauberentry| \textit{\empty}, wird der |\zaubername| auf der
% Gliederungsebene |subsubsection| zum Inhaltsverzeichnis
% hinzugefügt. Andernfalls geschieht dasselbe mit |\zauberentry|.
%    \begin{macrocode}
\phantomsection%
\ifdefstring{\zauberentry}{\empty}{%
        \addcontentsline{toc}{\midgard@zauber@tocLevel}{\zaubername}
        }{%
        \addcontentsline{toc}{\midgard@zauber@tocLevel}{\zauberentry}
        }
%    \end{macrocode}
% Zunächst wird der Name der Beschwörung gedruckt. Dann wird zunächst
% |\hfill| benutzt, um an das Ende der Zeile vorzurücken. 
%    \begin{macrocode}
\noindent\makebox[\linewidth]{\parbox[b]{0.85\linewidth}{%
    \unraggedright\expandnext{\midgard@htarget}{\zaubername}%
    \noindent \bfseries \large  \zaubername%
  }%
\hfill%
\ifdefstring{\zauberscroll}{true}%
           {}%
           {\imgMagisch}%
\hfill}\par%
\setlength{\parindent}{-1ex}%
\setlength{\leftskip}{1ex}%
\setlength{\rightskip}{0pt plus 0.2\linewidth}%
\setlength{\parskip}{0pt}%
\nopagebreak\par
%    \end{macrocode}
% Falls |\zauberAnmerkung| gesetzt ist, wird sie jetzt klein, fett und
% kursiv abgedruckt.
%    \begin{macrocode}
\ifdefstring{\zauberAnmerkung}{\empty}{}%
{{\par\footnotesize\bfseries\itshape\zauberAnmerkung\par}}%
%    \end{macrocode}
% Falls |\zauberUntertitel| gesetzt ist, wird sie jetzt klein, fett und
% kursiv abgedruckt.
%    \begin{macrocode}
\ifdefstring{\zauberUntertitel}{\empty}{}%
{{\nopagebreak\par\itshape\zauberUntertitel\par}}%
%    \end{macrocode}
% Wir unterdrücken einen Zeilenumbruch unmittelbar nach der
% Titelzeile. 
%    \begin{macrocode}
\nopagebreak\par%
%    \end{macrocode}
% Die Zauberart ist in jedem Fall \glqq Beschwörung\grqq. Ist die Stufe
% \textit{GM}, so wird der Text \glqq der Großen Magie\grqq{}
% nachgestellt, andernfalls \glqq der Stufe |\zauberstufe|\grqq{}
%    \begin{macrocode}
\vspace{\midgardskip}\par
Beschwörung~%
\ifdefstring{\zauberstufe}{GM}%
           {der Großen Magie}%
           {der Stufe \zauberstufe}%
\par
%    \end{macrocode}
% Anstelle des Zaubermaterials listen wir den Köder und das
% Beschwörungsfeld auf.
%    \begin{macrocode}
\def\beschwoerungsfeldtotal{}%
\ifdefstring{\beschwoerungsfeld}{\empty}%
           {}%
           {\appto\beschwoerungsfeldtotal{\beschwoerungsfeld}%
             \ifdefstring{\beschwoerungsfeldpreis}{\empty}%
                        {}%
                        {\appto\beschwoerungsfeldtotal{~(\beschwoerungsfeldpreis)}}%
           }%
\beschwoerungsfeldtotal\par%
\def\beschwoerungskoedertotal{}%
\ifdefstring{\beschwoerungskoeder}{\empty}%
            {}%
            {\appto\beschwoerungskoedertotal{\textit{\beschwoerungskoeder}}%
              \ifdefstring{\beschwoerungskoederpreis}{\empty}%
                          {}%
                          {\appto\beschwoerungskoedertotal{~(\beschwoerungskoederpreis)}}%
            }%
\ifdefstring{\beschwoerungskoeder}{\empty}%
            {}% 
            {\textbf{Köder:}~\beschwoerungskoedertotal\par}%
%    \end{macrocode}
% Gegebenenfalls drucken wir auch die Zauberformeln, falls dies
% gewünscht ist.
%    \begin{macrocode}
\vspace{\midgardskip}
\ifdefstring{\showformula}{\empty}{}{
  \beschwoerungprimprozess~%
  \imgModus~%
  \beschwoerungprimagens~%
  \imgObjectum~%
  \beschwoerungprimreagens\par%
  \beschwoerungsecprozess~%
  \imgModus~%
  \beschwoerungsecagens~%
  \imgObjectum~%
  \beschwoerungsecreagens\par%
  \vspace{\parskip}%
  \vspace{\midgardskip}%
}%
%    \end{macrocode}
% In einer kompaktifizierten |tabbing|-Umgebung drucken wir die
% wichtigen Daten jeweils, falls diese nicht auf \textit{hide}
% gesetzt sind.
%    \begin{macrocode}
\begin{spacelesstabbing}%
\textbf{Wirkungsbereich:}\hspace{1ex} \= anytext \kill
\ifdefstring{\zauberAP}{hide}{}{\textbf{AP-Verbrauch:} \> \zauberAP\\}%
\ifdefstring{\zauberGZ}{hide}{}{\textbf{Gefährdungszahl:} \> \zauberGZ\\}%
\ifdefstring{\zauberZd}{hide}{}{\textbf{Zauberdauer:} \> \zauberZd\\}%
\ifdefstring{\zauberUr}{hide}{}{\textbf{Ursprung:} \> \zauberUr\\}
\end{spacelesstabbing}\vspace{-1.0em}\par%
%    \end{macrocode}
% Wir unterdrücken wiederum etwas vertikalen Platz und drucken dann,
% jeweils falls gesetzt, die Lernkosten. Hierbei ist wichtig, dass wir
% den Fall \textit{niemand} gesondert behandeln, da in diesem Fall
% auch die Kosten selbst nicht abgedruckt werden sollen.
%    \begin{macrocode}
\vspace{1.0\midgardskip}\par
\ifdefstring{\zauberFP}{hide}{}{
  \ifdefstring{\zauberG}{niemand}%
  {niemand }%
  {\textbf{\themidgard@GrundKosten :} \nolinebreak \zauberG} -- %
  \ifdefstring{\zauberS}{niemand}%
  {niemand }%
  {\textbf{\themidgard@StandardKosten :} \nolinebreak \zauberS} -- %
  \ifdefstring{\zauberA}{niemand}%
  {niemand }%
  {\textbf{\themidgard@AusnahmeKosten :} \nolinebreak \zauberA}\par%
  \vspace{1.6\midgardskip}
}%
\expandafter\ifcsdef{\zauberKey exportvalues}{%
  \ifdatablist{\adddataline{\zauberKey}{\csname \zauberKey exportvalues\endcsname}}%
}{}%
}%
\endgroup%
%    \end{macrocode}
% Abschließend sorgen wir für einen ausreichenden Abstand und
% verhindern die Einrückung des auf den Datenblock folgenden Absatzes.
%    \begin{macrocode}
\@afterindentfalse%
\@afterheading%
}%
%    \end{macrocode}
% Schlussendlich definieren wir noch ein Kürzel für Anrufungen-, Namens- und Geisterbeschwörungen,
%    \begin{macrocode}
\newcommand*\individualbeschwoerung[2][]{%
  \beschwoerung[%
    secreagens=Luft,%
    showformula=true,%
    #1%
  ]{#2}%
}
%    \end{macrocode}
%\end{macro}
%
%\begin{macro}{\zaubermittel}
%  Für die Zaubermittel wählen wir ein ähnliches Vorgehen wie für
%  das Arkanums-Layout der Zauber. Zunächst definieren wir lokale
%  Kommandos.
%    \begin{macrocode}
\setzeGliederungsEbene{zaubermittel}{subsubsection}
\define@key{zaubermittel}{name}{\def\zaubermittelName{#1}}
\define@key{zaubermittel}{entry}{\def\zaubermittelEntry{#1}}
\define@key{zaubermittel}{voraussetzungen}{\def\zaubermittelVor{#1}}
\define@key{zaubermittel}{stufe}{\def\zaubermittelStufe{#1}}
\define@key{zaubermittel}{zeitaufwand}{\def\zaubermittelZeit{#1}}
\define@key{zaubermittel}{kosten}{\def\zaubermittelKosten{#1}}
\define@key{zaubermittel}{FP}{\def\zaubermittelFP{#1}}
\define@key{zaubermittel}{G}{\def\zaubermittelG{#1}}
\define@key{zaubermittel}{S}{\def\zaubermittelS{#1}}
\define@key{zaubermittel}{A}{\def\zaubermittelA{#1}}
\count@=0
\loop
  \advance\count@ 1
  \begingroup\edef\x{\endgroup
    \noexpand\define@key{zaubermittel}{FP\the\count@}[]{%
      \noexpand\@namedef{zaubermittelFP\the\count@}{####1}%
    }% end of \define@key
  }% end of \edef
  \x % execute \x
\ifnum\count@<6
\repeat
%    \end{macrocode}
% Wir definieren \textit{default}-Werte für die Variablen.
%    \begin{macrocode}
\presetkeys{zaubermittel}{voraussetzungen=keine,stufe=1,zeitaufwand=keine,
  FP=keine,G=niemand,S=niemand,A=niemand,entry=\empty,kosten=--,zeitaufwand=--,
  FP1=0,FP2=0,FP3=0,FP4=0,FP5=0,FP6=0,%
}{}
%    \end{macrocode}
%    \begin{macrocode}
\newcount\tmp@count@%
\newcommand*\zaubermittel[2][]{%
\setkeys{zaubermittel}{#1}{
\setlength{\parindent}{0pt}\par
%\setlength{\parskip}{0pt}\par
\ifdefstring{\zaubermittelFP}{keine}{}{

}%
\pagebreak[2]
\vspace{1em}
\phantomsection%
\ifdefstring{\zaubermittelEntry}{\empty}{%
        \addcontentsline{toc}{\midgard@zaubermittel@tocLevel}{\zaubermittelName}
        }{%
        \addcontentsline{toc}{\midgard@zaubermittel@tocLevel}{\zaubermittelEntry}
      }
\noindent{\bfseries\zaubermittelName} -- Stufe \zaubermittelStufe%
\ifdefstring{\zaubermittelVor}{keine}{}{\\\textit{\zaubermittelVor}}%
\par%
\ifdefstrequal{\zaubermittelZeit}{\zaubermittelKosten}{}{
  \vspace{0.5em}%
  \begin{spacelesstabbing}%
    \textbf{Zeitaufwand:}\hspace{1ex} \= anytext \kill
    \textbf{Zeitaufwand:} \> \zaubermittelZeit\\
    \textbf{Kosten:} \> \zaubermittelKosten
  \end{spacelesstabbing}%
}
\vspace{0.5em}%
\ifdefstring{\zaubermittelFP}{keine}{%
  \begin{enumerate}[align=left,itemsep=0pt,topsep=0pt,parsep=0pt]
  \tmp@count@=0%
  \loop%
  \advance\tmp@count@ 1%
  \ifdefstring{\csname zaubermittelFP\the\tmp@count@\endcsname}{0}{}{%
    \setcounter{midgard@StandardKosten}{\csname zaubermittelFP\the\tmp@count@\endcsname}%
    \setcounter{midgard@GrundKosten}{\value{midgard@StandardKosten}/2}%
    \setcounter{midgard@AusnahmeKosten}{\value{midgard@StandardKosten}*5}%
    \ifnum\value{midgard@StandardKosten}>0\relax%
    \item[Stufe \the\tmp@count@:] 
    \ifdefstring{\zaubermittelG}{niemand}%
                {niemand }%
                {\textbf{\themidgard@GrundKosten :} \nolinebreak \zaubermittelG} -- %
    \ifdefstring{\zaubermittelS}{niemand}%
                {niemand }%
                {\textbf{\themidgard@StandardKosten :} \nolinebreak \zaubermittelS} -- %
    \ifdefstring{\zaubermittelA}{niemand}%
                {niemand }%
                {\textbf{\themidgard@AusnahmeKosten :} \nolinebreak \zaubermittelA}%
    \fi%
  }%
  \ifnum\tmp@count@<6\repeat%
  \end{enumerate}
  \vspace{0.5em}%
}{
  \setcounter{midgard@StandardKosten}{\zaubermittelFP}%
  \setcounter{midgard@GrundKosten}{\value{midgard@StandardKosten}/2}%
  \setcounter{midgard@AusnahmeKosten}{\value{midgard@StandardKosten}*5}%
  \ifdefstring{\zaubermittelG}{niemand}%
  {niemand }%
  {\textbf{\themidgard@GrundKosten :} \nolinebreak \zaubermittelG} -- %
  \ifdefstring{\zaubermittelS}{niemand}%
  {niemand }%
  {\textbf{\themidgard@StandardKosten :} \nolinebreak \zaubermittelS} -- %
  \ifdefstring{\zaubermittelA}{niemand}%
  {niemand }%
  {\textbf{\themidgard@AusnahmeKosten :} \nolinebreak \zaubermittelA}%
  \par\vspace{0.5em}%
}
}%
\@afterindentfalse%
\@afterheading%
}%
%    \end{macrocode}
% \end{macro}
%
% \begin{macro}{\zauberwerkstatt}
%   Aus Gründen der Kompatiblität und Bequemlichkeit definieren wir
%   das Kommand |\zauberwerkstatt| als identisch mit |\zaubermittel|.
%    \begin{macrocode}
\let\zauberwerkstatt\zaubermittel
%    \end{macrocode}
% \end{macro}
%
%\begin{macro}{\artefakt}
% Auch Artefakte haben im Arkanum ein charakteristisches Layout,
% welches jedoch weniger komplex ist als beispielsweise die
% Zauberdatenblöcke. Die Herausforderung besteht hier darin,
% vernünftige Zeilenumbrüche zu erzeugen.
%    \begin{macrocode}
\setzeGliederungsEbene{artefakt}{subsubsection}
\define@key{artefakt}{name}            {\def\artefaktName{#1}}
\define@key{artefakt}{ABW}             {\def\artefaktABW{#1}}
\define@key{artefakt}{aura}            {\def\artefaktAura{#1}}
\define@key{artefakt}{typ}             {\def\artefaktTyp{#1}}
\define@key{artefakt}{magisch}         {\def\artefaktMagisch{#1}}
\define@key{artefakt}{mag}         {\def\artefaktMagisch{#1}}
\define@key{artefakt}{herstellbar}     {\def\artefaktHerstellbar{#1}}
\define@key{artefakt}{untertitel}      {\def\artefaktUntertitel{#1}}
\define@key{artefakt}{sprache}         {\def\artefaktSprache{#1}}
\define@key{artefakt}{farbe}           {\def\artefaktSprache{#1}}
\presetkeys{artefakt}{name=Artefakt,ABW=-,aura=-,herstellbar=false,typ=-,
  untertitel=-,sprache=-,magisch=true%
}{}
%    \end{macrocode}
% Von Spielerfiguren herstellbare Artefakten haben einen erweiterten
% Datenblock. Um die Implementierung und die Bedienung so einfach wie
% möglich zu halten, wird dieser zweite Datenblock von einem separaten
% Kommando erzeugt.
%    \begin{macrocode}
\define@key{artefaktherstellen}{zauber}{\def\herstellenVoraussetzung{#1}}
\define@key{artefaktherstellen}{stufe} {\def\herstellenStufe{#1}}
\define@key{artefaktherstellen}{ABW}   {\def\herstellenABW{#1}}
\define@key{artefaktherstellen}{EP}    {\def\herstellenEP{#1}}
\define@key{artefaktherstellen}{material}{\def\herstellenMaterial{#1}}
\presetkeys{artefaktherstellen}{zauber=-,stufe=-,ABW=-,EP=0}{}
%    \end{macrocode}
% Das |\artefakt|-Kommando erzeugt den Artefakt-Datenblock.
%    \begin{macrocode}
\newcommand\midgard@artefakt@customcode\relax%
\newcommand*{\artefakt}[2][]{
%    \end{macrocode}
% Erneut beginnen wir mit dem Vorschlag eines Seitenumbruches.
%    \begin{macrocode}
\par\pagebreak[2]%
\addvspace{1em}% plus 0.5em minus 0.5em}%
\setkeys{artefakt}{#1}{%
%    \end{macrocode}
% Wir fügen den Artefaktnamen dem Inhaltsverzeichnis hinzu.
%    \begin{macrocode}
\addcontentsline{toc}{\midgard@artefakt@tocLevel}{\artefaktName}
\begin{samepage}
  \setlength\parindent{0pt}
  \flushleft
      {\bfseries \large  \artefaktName}\ifdefstring{\artefaktHerstellbar}{false}{}{~\textdagger} \linebreak[0]%
      \ifdefstring{\artefaktTyp}{-}{}{\optendash\mbox{\artefaktTyp}}%
      \ifdefstring{\artefaktMagisch}{false}{\optendash nicht magisch}{%
        \ifdefstring{\artefaktABW}{-}{}{ \mbox{(\textsc{Abw}\,\artefaktABW)}}%
        \ifdefstring{\artefaktAura}{-}{}{ \optendash \mbox{\textbf{\textsc{Aura:}} \artefaktAura}}%
      } 
      \ifdefstring{\artefaktUntertitel}{-}{%
        \ifdefstring{\artefaktSprache}{-}{}{\newline\artefaktSprache}%
      }{\newline%
        \textit{\artefaktUntertitel}%
        \ifdefstring{\artefaktSprache}{-}{}{\optendash \artefaktSprache}%
      }
      \end{samepage}
      \par\nopagebreak
}\nopagebreak
%    \end{macrocode}
% Zuletzt verhindern wir Einrückung in der nachfolgenden Zeile.
%    \begin{macrocode}
\midgard@artefakt@customcode%
\@afterindentfalse%
\@afterheading%
}
%    \end{macrocode}
% Das |\artefaktherstellen|-Kommando erzeugt zweiten Datenblock für
% von Spielerfiguren herstellbare Artefakte.
%    \begin{macrocode}
\newcommand*{\artefaktherstellen}[2][]{
%    \end{macrocode}
% Erneut beginnen wir mit dem Vorschlag eines Seitenumbruches.
%    \begin{macrocode}
\begingroup\small%
\setkeys{artefaktherstellen}{#1}{%
\noindent%
\ifdefstring{\herstellenVoraussetzung}{-}{}{%
  \textit{\herstellenVoraussetzung}%
  \ifdefstring{\herstellenStufe}{-}{}{~(Stufe \herstellenStufe)} 
}
\herstellenEP\,EP
\ifdefstring{\herstellenABW}{-}{}{für \textsc{Abw}\,\herstellenABW}
\ifdefstring{\herstellenMaterial}{-}{}{%
\newline\noindent\textit{Material:}~\herstellenMaterial
}}\endgroup
\par
%    \end{macrocode}
% Zuletzt verhindern wir Einrückung in der nachfolgenden Zeile.
%    \begin{macrocode}
\@afterindentfalse%
\@afterheading%
}
%    \end{macrocode}
% \end{macro}
%
% \subsection{Bestiarium}
%
%    \begin{macrocode}
\define@key{best}{name}{\def\bestName{#1}}
\define@key{best}{typ}{\def\bestTyp{#1}}
\define@key{best}{modus}{\def\bestModus{#1}}
\define@key{best}{addendum}{\def\bestAdd{#1}}
\define@key{best}{variant}{\def\bestVariant{#1}}
\define@key{best}{grad}{\def\bestGrad{#1}}
\define@key{best}{mag} {\def\bestMag {#1}}
\define@key{best}{In}  {\def\bestIn  {#1}}
\define@key{best}{size}{\def\bestSize  {#1}}
\define@key{best}{EP}  {\def\bestEP  {#1}}
\define@key{best}{LP}  {\def\bestLP  {#1}}
\define@key{best}{energy}  {\def\bestEnergie  {#1}}
\define@key{best}{initiative}  {\def\bestInitiative  {#1}}
\define@key{best}{raufen}  {\def\bestRaufen  {#1}}
\define@key{best}{AP}  {\def\bestAP  {#1}}
\define@key{best}{Gw}  {\def\bestGw  {#1}}
\define@key{best}{St}  {\def\bestSt  {#1}}
\define@key{best}{B}   {\def\bestB   {#1}}
\define@key{best}{RK}  {\def\bestRK  {#1}}
\define@key{best}{MW}  {\def\bestMW  {#1}}
\define@key{best}{abwehr}{\def\bestAbwehr  {#1}}
\define@key{best}{resistenz}  {\def\bestResistenz  {#1}}
\define@key{best}{angriff}{\def\bestAngriff  {#1}}
\define@key{best}{bes}    {\def\bestBes  {#1}}
\define@key{best}{aura}    {\def\bestAura  {#1}}
\define@key{best}{vorkommen}    {\def\bestVork  {#1}}
\define@key{best}{zauber}             {\def\bestZauber {#1}}
\define@key{best}{zauberEW}           {\def\bestZauberEW {#1}}
\define@key{best}{zauberMore}         {\def\bestZauberMore {#1}}
\define@key{best}{zauberMMore}        {\def\bestZauberMMore {#1}}
\define@key{best}{zauberMoreEW}       {\def\bestZauberMoreEW {#1}}
\define@key{best}{zauberMMoreEW}      {\def\bestZauberMMoreEW {#1}}
\define@key{best}{zauberArtig}        {\def\bestZauberArtig {#1}}
\define@key{best}{zauberArtigEW}      {\def\bestZauberArtigEW {#1}}
\define@key{best}{zauberArtigMore}    {\def\bestZauberArtigMore {#1}}
\define@key{best}{zauberArtigMoreEW}  {\def\bestZauberArtigMoreEW {#1}}
\define@key{best}{zauberArtigMMore}   {\def\bestZauberArtigMMore {#1}}
\define@key{best}{zauberArtigMMoreEW} {\def\bestZauberArtigMMoreEW {#1}}
\define@key{best}{bestcategory}       {\def\bestKey{#1}}

\presetkeys{best}{ name=Wesen, grad=0, modus=\empty, typ=\empty,
  mag=false, aura=\empty, size=normal, addendum=\empty, variant=\empty,
  LP=\textstar, AP=\textstar, RK=\textstar R, EP=0, In=t01, 
  St=\textstar, Gw=01, B=\textstar, MW=\empty,
  abwehr=10, resistenz=10/10/10, angriff=keiner, zauber=\empty,
  zauberMore=\empty, zauberMMore=\empty, zauberArtig=\empty,
  zauberArtigMore=\empty, zauberArtigMMore=\empty, zauberEW=0,
  zauberMoreEW=0, zauberMMoreEW=0, zauberArtigEW=0,
  zauberArtigMoreEW=0, zauberArtigMMoreEW=0, bes=\empty, vorkommen=\empty,
  energy=0, initiative=\empty,raufen=\empty,bestcategory=best }{}
%    \end{macrocode}
%
% \begin{macro}{\bestiarium}
%  Wir benutzen das Paket |mdframed|, um eine Box darzustellen.
% 
%    \begin{macrocode}
\mdfdefinestyle{bestbox}{%
  skipabove=0.7em plus 0.3em minus 0.2em,%
  skipbelow=0.2em plus 0.5em minus 0.2em,%
  innertopmargin=0.5em,%
  innerbottommargin=0.6em,%
  innerleftmargin=1ex,%
  innerrightmargin=2ex}
%    \end{macrocode}
%
% Wir definieren das Kommando.
% 
%    \begin{macrocode}
\newcommand*\bestiarium[2][]{%
\setkeys{best}{#1}{
\par\setlength\parskip{0pt}\par
%    \end{macrocode}
% Das Kommando kennt die Optionen |noframe| und |float|,
% welche zum unterdrücken eines Rahmens führen bzw.~dazu, dass die
% umrahmte Box als Gleitobjekt generiert wird. Wird die Box nicht als
% Gleitobjekt generiert, wird ein Seitenumbruch vor der Box bevorzugt.
%    \begin{macrocode}
\ifdefstring{#2}{noframe}{%
  \pagebreak[2]
}{%
  \ifdefstring{#2}{float}{%
    \begin{figure}
  }{%
    \pagebreak[3]
  }
  \begin{mdframed}[style=bestbox]%
}%
%    \end{macrocode}
% Der Inhalt der Box soll einen hängenden Einzug haben.
%    \begin{macrocode}
\setlength{\leftskip}{1ex}
\setlength{\parindent}{-1ex}
\small%
\begin{spacelesstabbing}%
%    \end{macrocode}
% Die Spaltenbreiten des Tabellen-Layouts werden relativ zur Spaltenbreite gewählt.
%    \begin{macrocode}
\hspace{0.3\linewidth} \= \hspace{0.3\linewidth} \= \hspace{0.2\linewidth} \= anything\kill
%    \end{macrocode}
% Zunächst wird der name der Kreatur gesetzt.
%    \begin{macrocode}
\textbf{\bestName}%
%    \end{macrocode}
% Wurde die option |mag| verwendet, wird das Symbol \imgMagisch{} hinter den Namen gesetzt.
%    \begin{macrocode}
\ifdefstring{\bestMag}{true}%
           {~\imgMagisch}%
           {}% 
%    \end{macrocode}
% Ein Namenszusatz, der über |add| gesetzt wurde, folgt mit einem Komma abgetrennt.
%    \begin{macrocode}
\ifdefstring{\bestAdd}{\empty}%
           {}%
           {\textbf{,\,\bestAdd}}%
%    \end{macrocode}
% Eine Kreaturenvariante, die mit |variant| definiert wurde, folgt mit Komma getrennt und in normaler Wichtung.
%    \begin{macrocode}
\ifdefstring{\bestVariant}{\empty}%
           {}%
           {,\,\bestVariant}%
%    \end{macrocode}
% Der mit |mode| deklarierte Modus der Kreatur folgt etwas abgesetzt.
%    \begin{macrocode}
~%
\ifdefstring{\bestModus}{\empty}%
           {}
           {\bestModus~}% 
%    \end{macrocode}
% Die Klammer beinhaltet den Typ der Kreatur (sofern mit |type|
% definiert) sowie ihren mit |grad| gesetzten Grad und gegebenenfalls
% mit einem Gedankenstrich abgesetzt die über |size| definierte Größe.
%    \begin{macrocode}
(%
\ifdefstring{\bestTyp}{\empty}%
           {}%
           {\bestTyp~}%
Grad \bestGrad%
\ifdefstring{\bestSize}{normal}{}{ -- \bestSize}%
)%
%    \end{macrocode}
% Der Rest der Spalte beinhaltet die über |In| definierte Intelligenz.
%    \begin{macrocode}
 \> \> \> \ifdefstring{\bestIn}{hide}{~}{In: \bestIn}\\
%    \end{macrocode}
% In der nächsten Spalte folgen die Lebenspunkte |LP|, Ausdauerpunkte
% |AP|, der Moralwert |MW| und der EP-Multiplikator |EP| mit passender
% Beschriftung. Für Lebens- und Ausdauerpunkte werden die
% Beschriftungen fett gesetzt, für den Moralwert das \glqq
% +\grqq{}-Zeichen automatisch eingefügt.
%    \begin{macrocode}
\textbf{LP} \bestLP \> \textbf{AP} \bestAP \> \ifdefstring{\bestMW}{\empty}{}{MW+\bestMW}\> EP \bestEP\\
%    \end{macrocode}
% Es folgen Gewandtheit |Gw|, Stärke |St|, Bewegungsweite |B| und Rüstungsklasse |RK|, letztere in Fettdruck.
%    \begin{macrocode}
Gw \bestGw \> St \bestSt \> B\,\bestB \> \textbf{\bestRK}\\
%    \end{macrocode}
% Als letzte im Tabellenlayout definierte Zeile folgen Abwehr
% |abwehr|, Resistenz gegen Magie |resisenz| (mit automatisch
% hinzugefügten \glqq +\grqq{}-Zeichen) und die Energie |energie|,
% falls definiert.
%    \begin{macrocode}
Abwehr+\bestAbwehr \> Resistenz+\bestResistenz \ifdefstring{\bestEnergie}{0}{}{\> \textsc{E} \bestEnergie}
\end{spacelesstabbing}\par
%    \end{macrocode}
% Die Angriffstechniken |angriff| der Kreatur werden als entsprechend überschriebener Absatz mit hängendem Einzug gesetzt.
%    \begin{macrocode}
\textbf{Angriff:} \sloppy \bestAngriff%
%    \end{macrocode}
% Falls gewünscht wird der Eintrag für \textit{Raufen} automatisch aus dem Wert |raufen| generiert.
%    \begin{macrocode}
\ifdefstring{\bestRaufen}{\empty}{}{ -- Raufen+\bestRaufen}%
%    \end{macrocode}
% An dieser Stelle kann auch die Initiative |initiative| abgedruckt werden.
%    \begin{macrocode}
\ifdefstring{\bestInitiative}{\empty}{}{ -- \textsc{Initiative}+\bestInitiative}%
\par
%    \end{macrocode}
% Die Besonderheiten |bes| der Kreatur werden als entsprechend
% überschriebener Absatz mit hängendem Einzug gesetzt. Nötigenfalls
% wird ein Eintrag für die Aura |aura| generiert, wobei die
% Beschriftung in fetten Kapitälchen erfolgt.
%    \begin{macrocode}
\ifdefstring{\bestBes}{\empty}%
           {\ifdefstring{\bestAura}{\empty}{}{\textbf{\textsc{Aura:}} \bestAura\par}}%
           {\textbf{Bes.:} \bestBes \ifdefstring{\bestAura}{\empty}{}{ \optendash \textbf{\textsc{Aura:} \bestAura}}\par }%
%    \end{macrocode}
% Die Zauberfertigkeiten |zauber| der Kreatur werden als mit
% \textbf{Zaubern+|zauberEW|:} überschriebener Absatz mit hängendem
% Einzug gesetzt. Entsprechende Einträge werden auch generiert für
% |zauberMore| mit |zauberMoreEW|, |zauberMMore| mit |zauberMMoreEW|,
% |zauberArtig| mit |zauberArtigEW|, |zauberArtigMore| mit
% |zauberArtigMoreEW| und |zauberArtigMMore| mit |zauberArtigMMoreEW|.
%    \begin{macrocode}
\ifdefstring{\bestZauber}{\empty}%
           {}%
           {\textbf{Zaubern+\bestZauberEW:} \textit{\bestZauber}\par}%
\ifdefstring{\bestZauberMore}{\empty}%
           {}%
           {\textbf{Zaubern+\bestZauberMoreEW:} \textit{\bestZauberMore}\par}%
\ifdefstring{\bestZauberMMore}{\empty}%
           {}%
           {\textbf{Zaubern+\bestZauberMMoreEW:} \textit{\bestZauberMMore}\par}%
\ifdefstring{\bestZauberArtig}{\empty}%
           {}%
           {\textbf{\guillemotright Zaubern+\bestZauberArtigEW\guillemotleft:} \textit{\bestZauberArtig}\par}%
\ifdefstring{\bestZauberArtigMore}{\empty}%
           {}%
           {\textbf{\guillemotright Zaubern+\bestZauberArtigMoreEW\guillemotleft:} \textit{\bestZauberArtigMore}\par}%
\ifdefstring{\bestZauberArtigMMore}{\empty}%
           {}%
           {\textbf{\guillemotright Zaubern+\bestZauberArtigMMoreEW\guillemotleft:} \textit{\bestZauberArtigMMore}\par}%
%    \end{macrocode}
% Das Vorkommen der Kreatur letztlich wird in ähnlichem Stil mit hängendem Einzug erzeugt.           
%    \begin{macrocode}
\ifdefstring{\bestVork}{\empty}%
           {}%
           {\textbf{Vorkommen:} \bestVork}%
%    \end{macrocode}
% Auch beim schließen der Umgebungen wird ein Seitenumbruch vorgeschlagen.
%    \begin{macrocode}
\ifdefstring{#2}{noframe}{%
  \pagebreak[2]
}{%
  \end{mdframed}%
  \ifdefstring{#2}{float}{%
    \vspace{-1.5em}
    \end{figure}
  }{%
    \pagebreak[3]
  }
}%
%    \end{macrocode}
% Letztlich können auch beim Bestiariums-Kommate |\bestiarum| Werte nach belieben exportiert werden.
%    \begin{macrocode}
\expandafter\ifcsdef{\bestKey exportvalues}{%
  \ifdatablist{\adddataline{\bestKey}{\csname \bestKey exportvalues\endcsname}}%
}{}%
}%
\@afterindentfalse%
\@afterheading%
}%
%    \end{macrocode}
%\end{macro}
%\begin{macro}{\best}
% Auf analoge Weise definieren wir auch eine Kurzversion des des
% Bestiarums-Layouts, welche im Fließtext verwendet werden kann.
%    \begin{macrocode}
\newcommand*\best[2][]{%
\pagebreak[3]
\setkeys{best}{#1}{
\setlength{\parskip}{0pt}
\setlength{\topsep}{0pt}%
\setlength{\partopsep}{0pt}\par
\small
\vspace{0.5em plus 0.5em minus 0.2em}%
\begin{spacing}{0.5}%
\noindent \textbf{\bestName}%
\ifdefstring{\bestMag}{true}{~\imgMagisch}{}%
\ifdefstring{\bestAdd}{\empty}{}{\textbf{,\,\bestAdd}}%
\ifdefstring{\bestVariant}{\empty}{}{,\,\bestVariant}%
~\ifdefstring{\bestModus}{\empty}{}{\bestModus~}%
(%
\ifdefstring{\bestTyp}{\empty}{}{\bestTyp~}%
Grad \bestGrad%
\ifdefstring{\bestSize}{normal}{}{ -- \bestSize}%
)\par%
\noindent\bestLP~LP, \bestAP~AP -- \bestRK~-- St~\bestSt, Gw~\bestGw, B\,\bestB \ifdefstring{\bestEnergie}{0}{}{, \textsc{E}\,\bestEnergie}\par%
\setlength{\leftskip}{1ex}
\setlength{\parindent}{-1ex}
\textbf{Angriff:} \bestAngriff%
\ifdefstring{\bestRaufen}{\empty}{}{ -- Raufen+\bestRaufen}%
\ifdefstring{\bestInitiative}{\empty}{}{ -- \textsc{Initiative}+\bestInitiative}%
~-- \textsc{Abwehr}+\bestAbwehr, \textsc{Resistenz}+\bestResistenz%
\ifdefstring{\bestBes}{\empty}{%
  \ifdefstring{\bestAura}{\empty}{}{\par\textsc{Aura:} \bestAura}%
}{%
  \par\textbf{Bes.:} \bestBes%
  \ifdefstring{\bestAura}{\empty}{}{ -- \textsc{Aura:} \bestAura}%
}%
\ifdefstring{\bestZauber}{\empty}{}{\par\textbf{Zaubern+\bestZauberEW:} \textit{\bestZauber}}%
\ifdefstring{\bestZauberMore}{\empty}{}{\par\textbf{Zaubern+\bestZauberMoreEW:} \textit{\bestZauberMore}}%
\ifdefstring{\bestZauberMMore}{\empty}{}{\par\textbf{Zaubern+\bestZauberMMoreEW:} \textit{\bestZauberMMore}}%
\ifdefstring{\bestZauberArtig}{\empty}{}{\par\textbf{\guillemotright Zaubern+\bestZauberArtigEW\guillemotleft:} \textit{\bestZauberArtig}}%
\ifdefstring{\bestZauberArtigMore}{\empty}{}{\par\textbf{\guillemotright Zaubern+\bestZauberArtigMoreEW\guillemotleft:} \textit{\bestZauberArtigMore}}%
\ifdefstring{\bestZauberArtigMMore}{\empty}{}{\par\textbf{\guillemotright Zaubern+\bestZauberArtigMMoreEW\guillemotleft:} \textit{\bestZauberArtigMMore}}%
\end{spacing}}\par%
\vspace{0.5em plus 0.5em minus 0.2em}%
\pagebreak[3]
\@afterindentfalse%
\@afterheading%
}%
%    \end{macrocode}
% \end{macro}
%
% \begin{macro}{\person}
%    \begin{macrocode}
\define@key{person}{name}{\def\personName{#1}}%
\define@key{person}{grad}{\def\personGrad{#1}}%
\define@key{person}{typ}{\def\personTyp{#1}}%
\define@key{person}{klasse}{\def\personKlasse{#1}}%
\define@key{person}{alter}{\def\personAlter{#1}}%
\define@key{person}{LP}{\def\personLP{#1}}%
\define@key{person}{AP}{\def\personAP{#1}}%
\define@key{person}{RK}{\def\personRK{#1}}%
\define@key{person}{B}{\def\personB{#1}}%
\define@key{person}{statur}{\def\personStatur{#1}}%
\define@key{person}{gewicht}{\def\personGewicht{#1}}%
\define@key{person}{groesse}{\def\personGroesse{#1}}%
\define@key{person}{aussehen}{\def\personAussehen{#1}}%
\define@key{person}{angriff}{\def\personAngriff{#1}}%
\define@key{person}{raufen}{\def\personRaufen{#1}}%
\define@key{person}{abwehr}{\def\personAbwehr{#1}}%
\define@key{person}{energie}{\def\personEnergie{#1}}%
\define@key{person}{initiative}{\def\personInitiative{#1}}%
\define@key{person}{resistenz}{\def\personResistenz{#1}}%
\define@key{person}{beschreibung}{\def\personBeschreibung{#1}}%
\define@key{person}{schicht}{\def\personSchicht{#1}}%
\define@key{person}{stand}{\def\personSchicht{#1}}%
\define@key{person}{herkunft}{\def\personHerkunft{#1}}%
\define@key{person}{glaube}{\def\personGlaube{#1}}%
\define@key{person}{fertigkeiten}{\def\personFertigkeiten{#1}}%
\define@key{person}{zauber}{\def\personZauber{#1}}%
\define@key{person}{zauberEW}{\def\personZauberEW{#1}}%
\define@key{person}{attribute}{\def\personAttribute{#1}}%
\define@key{person}{besitz}{\def\personBesitz{#1}}%
\define@key{person}{bes}{\def\personBesonderheiten{#1}}%
\presetkeys{person}{name=Person,grad=1,alter=,statur=,gewicht=,aussehen=,beschreibung=,schicht=,glaube=,herkunft=,fertigkeiten=,zauber=,zauberEW=,attribute=,typ=Normaler Mensch,besitz=,groesse=,klasse=,LP=,AP=,RK=,B=24,energie=0,raufen=,initiative=,bes=,angriff=}{}%
%    \end{macrocode}
%    \begin{macrocode}
\newcommand*\person[2][]{%
\pagebreak[3]
\setkeys{person}{#1}{
\setlength{\parskip}{0pt}
\setlength{\topsep}{0pt}%
\setlength{\partopsep}{0pt}\par
\small
\vspace{0.5em plus 0.5em minus 0.2em}%
\begin{spacing}{0.5}%
\noindent\textbf{\personName}\optendash\personTyp\hfill\mbox{}\linebreak[0]\mbox{}\nolinebreak\hfill\textbf{\ifdefstring{\personKlasse}{}{}{\mbox{\personKlasse{}}~}Gr.\,\personGrad}%
\par%
\setlength{\leftskip}{1ex}
\setlength{\parindent}{-1ex}
\ifdefstring{\personAttribute}{}{}{\personAttribute \optendash}
\ifdefstring{\personAlter}{}{}{\personAlter, }%
\ifdefstring{\personGroesse}{}{}{\personGroesse, }%
\ifdefstring{\personGewicht}{}{}{\personGewicht, }%
\ifdefstring{\personStatur}{}{}{\textit{\personStatur} \optendash}%
\ifdefstring{\personHerkunft}{}{}{\personHerkunft, }%
\ifdefstring{\personSchicht}{}{}{\personSchicht, }%
\ifdefstring{\personGlaube}{}{}{\textit{\personGlaube} \optendash}%
\ifdefstring{\personAussehen}{}{}{\personAussehen \optendash}%
\ifdefstring{\personBeschreibung}{}{}{\personBeschreibung}\par%
\ifdefstring{\personLP}{}{}{%
  \personLP~LP, \personAP~AP -- \personRK~-- B\,\personB \ifdefstring{\personEnergie}{0}{}{, \textsc{E}\,\personEnergie}\par%
  \ifdefstring{\personAngriff}{}{}{%
    \textbf{Angriff:} \personAngriff%
    \ifdefstring{\personRaufen}{}{}{ -- Raufen+\personRaufen}%
    \ifdefstring{\personInitiative}{}{}{ -- \textsc{Initiative}+\personInitiative}%
    ~-- \textsc{Abwehr}+\personAbwehr, \textsc{Resistenz}+\personResistenz%
  }
}
\ifdefstring{\personBesonderheiten}{\empty}{%
  \ifdefstring{\personAura}{\empty}{}{\par\textsc{Aura:} \personAura}%
}{}%
\ifdefstring{\personFertigkeiten}{}{}{\textit{\personFertigkeiten}\par}%
\ifdefstring{\personZauber}{}%
            {}%
            {\textbf{Zaubern+\personZauberEW:} \textit{\personZauber}\par}%
\ifdefstring{\personBesonderheiten}{}{}{
  \par\textbf{Bes.:} \personBesonderheiten%
  \ifdefstring{\personAura}{\empty}{}{-- \textsc{Aura:} \personAura}%
}%
\ifdefstring{\personBesitz}{}{}{\textbf{Besitz:} \personBesitz}\par%
\end{spacing}%
\bigskip
}
%    \end{macrocode}
% Zuletzt verhindern wir Einrückung in der nachfolgenden Zeile.
%    \begin{macrocode}
\@afterindentfalse%
\@afterheading%
}
%    \end{macrocode}
% \end{macro}
% \PrintChanges
% \PrintIndex
% \Finale
% \endinput
